\documentclass[11pt,openany]{book}

\usepackage[utf8]{inputenc}
\usepackage[T1]{fontenc}
\usepackage[english]{babel}
\usepackage{amsmath,amssymb,amsthm}
\usepackage{booktabs}
\usepackage{geometry}
\usepackage{hyperref}
\usepackage{listings}
\usepackage{longtable}
\usepackage{microtype}
\usepackage{url}
\usepackage{float}
\usepackage{tikz}
\usetikzlibrary{arrows.meta}
% Shared geometry and monochrome styles keep both editions visually consistent.
\tikzset{
  qsb node/.style={circle,draw,minimum size=8mm,inner sep=2pt,font=\small},
  qsb activity/.style={rectangle,draw,rounded corners=2pt,align=center,
    minimum width=16mm,minimum height=12mm,font=\small},
  qsb arc/.style={-{Stealth[length=2mm]},line width=0.6pt},
  qsb selected/.style={qsb arc,line width=1.4pt},
  qsb alternative/.style={qsb arc,dashed},
  qsb label/.style={font=\small,fill=white,inner sep=2pt},
  qsb terminal/.style={rectangle,draw,align=center,font=\small,inner sep=4pt}
}


\geometry{a4paper,margin=2.7cm}
\setlength{\emergencystretch}{2em}
\hypersetup{
  colorlinks=true,
  linkcolor=black,
  citecolor=black,
  urlcolor=blue,
  pdftitle={Operations Research with QSB: An Academic and Practical Manual},
  pdfauthor={QSB Project}
}


\theoremstyle{definition}
\newtheorem{definicion}{Definition}[chapter]
\newtheorem{ejemplo}{Example}[chapter]
\newtheorem{observacion}{Remark}[chapter]

\newcommand{\qsb}{\textnormal{\textsc{QSB}}}
\newcommand{\R}{\mathbb{R}}
\newcommand{\Z}{\mathbb{Z}}

\lstdefinestyle{qsbjson}{
  basicstyle=\ttfamily\small,
  breaklines=true,
  columns=fullflexible,
  frame=single,
  showstringspaces=false,
  keepspaces=true
}
\lstset{style=qsbjson}

\begin{document}

\frontmatter

\begin{titlepage}
  \centering
  \vspace*{2cm}
  {\Huge\bfseries Operations Research with \qsb\par}
  \vspace{0.5cm}
  {\Large\bfseries An Academic and Practical Manual\par}
  \vspace{1.5cm}
  {\Large Using the application, modeling problems, and interpreting results\par}
  \vspace{2cm}
  {\large QSB Project\par}
  \vspace{0.4cm}
  {\large Editors: the QSB academic and development community\par}
  \vfill
  {\large Verified working edition\par}
  {\large October 11, 2026\par}
\end{titlepage}

\chapter*{Editorial note}
\addcontentsline{toc}{chapter}{Editorial note}
This manual accompanies the QSB macOS application and command-line tool. It
presents operations research fundamentals, original reproducible examples, and
instructions for formulating, validating, solving, and interpreting models.
Its coverage matches the project's documented state on October 11, 2026. The
roadmap remains the public source of truth for software status.

The examples in \path{manual/examples/} were created for this project. They do
not require WinQSB files or reproduce its manual. The compatibility sections
explain how to import your own files without redistributing proprietary material.

\section*{How to read this manual}
Start with the application guide and the linear programming example. Then turn
to the chapter for the model family you need. Each chapter separates model
assumptions, the procedure in QSB, results, and method limitations. A numerical
solution does not by itself certify that the formulation accurately represents
the real problem.

Run commands from the repository root. JSON numbers use a decimal point; this
English edition also uses decimal points in mathematical text. Results are
rounded for readability. Verification compares numbers using tolerances rather
than textual equality of floating-point values.

\section*{Reproducibility}
Input files are part of the manual's sources. The verifier runs the associated
commands and checks the expected results. The manual distinguishes native
educational solvers from optional external engines, and does not attribute
guarantees of optimality to heuristics or local methods.

Both language editions share the same JSON examples and bibliography. Names
stored in those examples may be Spanish, such as \texttt{Lanzar} or
\texttt{Operativa}. They are retained as exact identifiers in this edition;
English explanations do not rename the underlying data. Application labels,
commands, JSON keys, units, numerical values, and solver guarantees retain the
same meaning in both editions.

\tableofcontents
\listoftables

\mainmatter

\part{Using QSB and modeling fundamentals}
\chapter{Practical application guide}
\label{chap:guia}

\section{A reproducible first walkthrough}
QSB provides a macOS application and a command-line tool called \texttt{qsb}.
Both share the models, validators, and solvers in \texttt{QSBCore}. The current
interface uses English labels; this manual retains those labels so they can be
located on screen.

To begin, open \path{manual/examples/lp_sample.json} using
\textbf{Model > Open Model...}. The model maximizes $50X_1+60X_2$ subject to
$2X_1+3X_2\leq180$, $3X_1+2X_2\leq150$, and nonnegative variables.
\begin{enumerate}
 \item In \textbf{Definition}, check names, coefficients, relations, bounds,
 and variable types. The variables here are continuous.
 \item Click \textbf{Validate} and review \textbf{Validation}. A valid model
 is a consistent input; it is not yet a solution.
 \item In \textbf{Run}, select \textbf{QSB Native} and click the \textbf{Run}
 button; you can also use \textbf{Solve} in the toolbar or
 \textbf{Solve > Run Current Model} in the menu.
 \item In \textbf{Summary}, check $X_1=18$, $X_2=48$, and objective 3780.
 Both constraints hold with equality. Also review the solution status and
 metadata before interpreting the number as an optimum.
 \item Export the model and solution separately. Reopen the exported model
 to preserve a reproducible input.
\end{enumerate}
The \textbf{Samples} menu provides built-in examples, but these may differ from
the manual's files. To reproduce a cited result, use the file specified in its
chapter.

\section{Creating, importing, and editing}
\textbf{New Model...} opens the family selector and creates an initial model.
\textbf{Open Model...} accepts normalized JSON and recognized legacy formats.
The family determines which controls appear. Native editors let you work with
names, tables, and typed fields; other variants use a structured form and JSON
as an interchange surface.

The \textbf{Definition} panel defines the problem; \textbf{Overview} summarizes
its family and state. \textbf{JSON} lets you inspect its representation. When
\textbf{Apply JSON to ... Editor} is available, explicitly apply the edited text
to replace the native draft. Writing JSON does not mean the changes have
already been incorporated into the table. Review the diagnostics if applying the JSON fails and preserve the text you need to repair.

Some tables offer \textbf{Copy Table} and \textbf{Paste Table} to exchange ranges
with spreadsheets. Check headers, column order, separators, units, and names
after pasting. Malformed numeric fields are not silently converted to zeros.
In editors that flag incomplete imported dimensions, review the cells and
shown values before accepting \textbf{Use Displayed Dimensions}.

\section{Validating, running, and interpreting states}
The application distinguishes empty, editing, validating, invalid, and valid
model states. Execution has separate states: \textbf{Not run},
\textbf{Validating}, \textbf{Solving}, \textbf{Solved}, and \textbf{Failed}.
Read a Solved status alongside the mathematical condition reported by the
family: a method may return an approximation, local solution, or heuristic
policy rather than a global optimum.

\textbf{Validation} displays model errors and warnings. Correct errors before
solving; a warning requires judgment even when it allows you to proceed.
\textbf{Run Details} lets you review the backend and execution information.
When a long run displays progress, wait for completion before repeating Run.
There is no universal cancellation control; some editors' asynchronous paths
prevent an old response from replacing a new draft. This does not cancel the
background computation.

\begin{description}
 \item[QSB Native] Deterministic educational methods. Exactness and limits
 depend on the family and are explained in the corresponding chapters.
 \item[Validate only] Checks the input and produces diagnostics; it does not
 compute decisions or optimal values.
 \item[External solver] Integrations available only for certain families.
 Availability depends on configuration and the installed solver. The presence
 of this option does not imply universal external support.
\end{description}
For LP/ILP and QP, consult the HiGHS conditions in the optimization chapter.
Other families may report that the external backend is unavailable. Switch to
QSB Native or Validate only according to the work you want to do.

\section{Saving models and results}
The toolbar's \textbf{Export} button prepares JSON from the current draft. If a
native field cannot be converted into the model, QSB displays validation rather
than exporting old data. The \textbf{Model > Export Model JSON...} menu opens
export of the model text; to ensure recent fields are synchronized, first use
Validate or the toolbar's Export button. \textbf{Export Solution JSON...} saves
the existing result and is enabled when solution content is present.

These are distinct documents: the model describes the problem to solve, and the
solution records what a run returned. Output formats vary by family; some
include the model and others do not. Always preserve both, together with the
command, backend, software version, and study assumptions. Do not use a solution
document as a model input.

\section{Keyboard controls and alternatives to charts}
The application defines the following shortcuts:
\begin{center}
\begin{tabular}{ll}
\toprule
Action & Keys\\
\midrule
New Model & Cmd+N\\
Open Model & Cmd+O\\
Export Model JSON & Cmd+S\\
Export Solution JSON & Cmd+Shift+S\\
Validate Model & Cmd+Shift+V\\
Run Current Model & Cmd+R\\
Solve ILP & Cmd+I\\
\bottomrule
\end{tabular}
\end{center}
Buttons are enabled according to family and state. Solve ILP and Solve LP
Relaxation are specific to linear/integer programming; the relaxation does not
enforce integrality. Some editors use Cmd+Shift+N to add rows. Use Tab to move
through native controls and the menus to locate actions; focus behavior depends
on the editor.

Tables and JSON let you inspect data without relying on chart colors. Native
results include accessible labels and values in several families, but this
alone does not establish a complete VoiceOver review. Preserve the numerical
tables for reproducible analysis. The macOS window has a minimum size of 980 by
620 points; use scrollable areas and close the Inspector when you need more width.

\section{The terminal and normalized JSON}
From the repository root, SwiftPM builds and runs the tool:
\begin{lstlisting}
swift run qsb
swift run qsb validate manual/examples/lp_sample.json
swift run qsb solve-json manual/examples/lp_sample.json --backend native
\end{lstlisting}
Without arguments, \texttt{qsb} displays the command catalog. After
\texttt{swift build}, you can run \texttt{.build/debug/qsb} directly. To save
output and separate errors:
\begin{lstlisting}
.build/debug/qsb solve-json manual/examples/lp_sample.json \
  --backend native > /tmp/qsb-lp-solution.json 2> /tmp/qsb-lp-errors.txt
\end{lstlisting}
JSON uses decimal points, double-quoted strings, and \texttt{null} for optional
fields that allow it. It does not include comments. Each family has its own
schema: LP uses a direct object, while many other families use a wrapper with
\texttt{kind} and \texttt{model}. Copy the structure of the family's example,
keeping names, vectors, and matrices aligned.

\texttt{solve-json} detects the family of a normalized input.
\texttt{validate} accepts normalized or legacy inputs. \texttt{solve} targets
legacy models; do not interchange commands indiscriminately. The chapters
provide family-specific commands, such as \texttt{solve-dp-json} or
\texttt{solve-decision-json}. \texttt{validate-json} is specific to LP/ILP,
rather than a universal validator.

To import your own legacy input, the general procedure is:
\begin{lstlisting}
qsb import-legacy-json /ruta/al/modelo > /tmp/qsb-model.json
qsb validate /tmp/qsb-model.json
qsb solve-json /tmp/qsb-model.json --backend native
\end{lstlisting}
Here \texttt{qsb} assumes that the executable is on PATH; you can replace it
with \texttt{swift run qsb}. The original files must be locally accessible and
used in accordance with their permissions. The manual's examples are
independent of proprietary material.

\section{Troubleshooting}
\begin{description}
 \item[Unreadable JSON] Check syntax, key names, and types; a decoding error
 may occur before structured diagnostics can be generated.
 \item[Invalid model] Consult \texttt{isValid}, severity, code, message, and
 diagnostic path when included in the output. Correct the input and validate again.
 \item[No feasible solution] The structure may be valid even though constraints
 or demand cannot be satisfied. Review bounds and balances.
 \item[Unavailable backend] Check which families support it and its
 configuration. Validate only does not provide an alternative solver result.
 \item[Unexpected result] Review units, objective sign, state order,
 constraints, options, and exactness. Compare with simple calculations before
 changing algorithms.
\end{description}
The CLI exits successfully when it completes a command and unsuccessfully when
an error is thrown, including incorrect usage. Some validations return a report
with \texttt{isValid: false} normally; the exit code alone does not prove that
the model is valid. An automated workflow must also inspect the validation
document and result status. The absence of syntax errors does not demonstrate
that the model accurately represents the real system.

\chapter{Methodological foundations}
\label{chap:preliminares}

\section{Purpose of the manual}
Operations research builds and analyzes quantitative models to support decisions
under constraints. This manual combines foundations with practical use: each
example turns a problem statement into QSB data and explains what the result
means. The problems are original teaching examples, with normalized inputs
available in \texttt{manual/examples/}; they do not require proprietary manuals
or files from earlier programs.

You can start with the application guide and a linear programming example, or
go directly to the chapter for your problem family. Before running a model,
read its assumptions: two models using the same words may represent different
systems. Do not confuse deterministic inventory with random demand, an acyclic
route with a general network, or a local optimum with a global one.

\section{The modeling cycle}
A rigorous study records:
\begin{enumerate}
 \item the real system, time horizon, and decision to be supported;
 \item controllable variables, parameters, and units;
 \item assumptions, uncertainty, and the scope of the data;
 \item objective, constraints, and initial conditions;
 \item solution method, backend, and execution options;
 \item feasibility checks and interpretation;
 \item sensitivity, risks, and limits of application.
\end{enumerate}
Write down the unit of every parameter. If demand is expressed per year and
lead time in days, convert them to a common basis before calculating. A number
without a unit can produce a mathematically consistent solution that is
operationally incorrect.

\section{Validity and exactness}
Validation checks structure, dimensions, fields, and family-specific rules.
It guarantees neither feasibility nor appropriate assumptions. A solver may
prove optimality for a model, return a heuristic solution, or stop at a limit.
Consult the status, algorithm, exactness, and execution notes.

In the metadata, \texttt{exact} characterizes the procedure for the supported
problem; it does not imply symbolic arithmetic or infinite precision.
\texttt{fixtureScale} identifies an educational method intended for small cases.
Numerical solution uses tolerances; compare results within a margin consistent
with their scale rather than requiring decimal strings to match exactly.
Repeatability does not prove validity either: a formulation error can be
perfectly deterministic.

\section{Computational traceability}
To reproduce an analysis, preserve the input JSON, QSB version, command, backend,
options, and result. Also record any transformations from a spreadsheet. For
external data, include the source, license, date, and unit of analysis. A
solution document does not always contain the complete model; save it separately.

The commands in this manual run from the repository root. \texttt{swift run qsb}
compiles if necessary; \texttt{.build/debug/qsb} uses the built executable. Small
examples let you check results by hand before studying larger cases. Changing
an input requires validation and solution again, rather than merely editing a
presentation table.

\begin{observacion}
The computational result expresses the consequences of the encoded data and
assumptions. The analyst must justify those assumptions, compare the solution
with the real system, and communicate its limitations to the decision maker.
\end{observacion}

\section{Comparison and sensitivity}
Distinguish verification (the program computes what was specified) from model
validation (what was specified adequately represents the system). Check balances
and constraints using the returned values. Then vary demand, costs, or
probabilities while keeping their distributions and units consistent. A change
in recommendation after a small perturbation may reveal a fragile decision; a
tie may require secondary criteria outside the model. Do not extrapolate a
conclusion beyond the range analyzed.


\part{Optimization and decision making}
\chapter{Linear programming}
\label{chap:pl}

\section{What a linear model represents}
Linear programming (LP) allocates scarce resources among activities when
both the objective and the constraints are linear
\cite{dantzig1963linear,bertsimas1997introduction}. A variable may represent
units produced, hours assigned, or quantities transported. Its meaning
and unit must be established before writing coefficients.

\begin{definicion}
A continuous linear model has the form
\[
 \operatorname{opt}_{x\in\R^n} c^\top x,\qquad
 a_i^\top x\ \{\leq,=,\geq\}\ b_i\quad(i=1,\ldots,m),\qquad
 \ell_j\leq x_j\leq u_j.
\]
The operator denotes minimization or maximization. A bound may be absent;
a free variable admits both positive and negative values.
\end{definicion}

Linearity assumes proportionality and additivity: doubling an activity
doubles its resource use and contribution, and contributions add without
interaction effects. A continuous formulation allows divisibility. The data
are also deterministic over the modeled horizon. Discounts, start-ups,
indivisible batches, or variable yields require a different formulation. Validating
a file checks its structure and mathematical rules; it does not establish that
these assumptions accurately represent the organization.

\section{Example: product mix}
A plant manufactures two divisible products, \(X_1\) and \(X_2\). Their unit
contributions are 50 and 60 monetary units. Each unit uses two resources as
shown in the following table; capacities cover the same period as production.
\begin{table}[htbp]
\centering
\caption{Product mix data}
\begin{tabular}{lrrr}
\toprule
Resource & Per unit of \(X_1\) & Per unit of \(X_2\) & Capacity \\
\midrule
Resource 1 & 2 & 3 & 180 \\
Resource 2 & 3 & 2 & 150 \\
\bottomrule
\end{tabular}
\end{table}
The model is
\[
\begin{aligned}
 \max\quad &50X_1+60X_2\\
 \text{subject to}\quad &2X_1+3X_2\leq180,\\
 &3X_1+2X_2\leq150,\\
 &X_1,X_2\geq0.
\end{aligned}
\]
Capacities are hard constraints: exceeding them is not allowed. Minimum
demand requirements and fixed costs are not included; attributing them to this
result would change the problem.

\subsection{Reproducible input and execution}
The complete model is in \path{manual/examples/lp_sample.json}.
\par\noindent The keys
\texttt{sense}, \texttt{variableNames}, \texttt{objectiveCoefficients}, and
\texttt{constraints} define the objective and rows. Each row includes a name,
coefficients in variable order, a relation, and \texttt{rhs}. Bounds are
\texttt{lowerBounds} and \texttt{upperBounds}; \texttt{null} in an upper bound
means that there is no explicit maximum. The example uses
\texttt{continuous} types. Run from the repository root:
\begin{lstlisting}
swift run qsb solve-json manual/examples/lp_sample.json --backend native
\end{lstlisting}
This specific command returns \texttt{objectiveValue} and the
\texttt{variableValues} dictionary. It does not include a sensitivity report.

On macOS, open the normalized JSON in the
\texttt{Linear / Integer Programming} family. The \texttt{Definition} view
shows \texttt{Model}, \texttt{Variables and objective}, and the constraints.
To enter it manually, select \texttt{Maximize}, create two variables with
coefficients 50 and 60, type \texttt{Continuous}, a lower bound of zero, and no
upper bound; then enter the two rows with relation \(\leq\). Use
\texttt{Validate} and correct the diagnostics before pressing \texttt{Run LP}.
The editor shortcuts are Command--Shift--V to validate and Command--R to
solve. \texttt{Copy Table} and \texttt{Paste Table} transfer tables to and from
a spreadsheet. Preserve the column order and headers;
review the model after pasting. In narrow windows, the editor may require
horizontal scrolling.

\subsection{Result and interpretation}
\begin{table}[htbp]
\centering
\caption{Linear optimum verified with QSB}
\begin{tabular}{lr}
\toprule
Quantity & Value \\
\midrule
\(X_1\) & 18 \\
\(X_2\) & 48 \\
Total contribution & 3780 \\
\bottomrule
\end{tabular}
\end{table}
The solution view shows the objective, variables, and constraint details.
Check feasibility numerically:
\[
 2(18)+3(48)=180,\qquad3(18)+2(48)=150.
\]
Both resources are fully used; their slacks \(b_i-a_i^\top x\) are zero.
An active constraint holds as an equality at the optimum; a positive slack
indicates unused capacity. This interpretation depends on the relation:
for \(a_i^\top x\geq b_i\), the surplus is \(a_i^\top x-b_i\).

The feasible region is a polygon with vertices \((0,0)\), \((50,0)\),
\((18,48)\), and \((0,60)\). Their objective values are 0, 2500, 3780, and 3600.
This independent geometric check confirms the result for the two-variable
problem. In higher dimensions, enumerating vertices by hand is impractical.

\section{Duality and sensitivity: what can be inferred}
For this maximization problem with constraints \(Ax\leq b\), \(x\geq0\), the dual is
\[
 \min\ 180y_1+150y_2,\qquad
 2y_1+3y_2\geq50,\quad3y_1+2y_2\geq60,\quad y_1,y_2\geq0.
\]
Since both production variables are positive at the optimum, both dual
constraints hold as equalities. They yield
\(y_1=16\), \(y_2=6\), and \(180(16)+150(6)=3780\). The dual values
certify an upper bound equal to the achieved objective.

As long as this basis remains optimal, one additional unit of resource 1
is worth 16 monetary units, and one of resource 2 is worth 6. These are local
marginal values, not prices valid for every expansion. For example, changing
only the first capacity to \(180+\delta\) moves the intersection to
\[
 X_1=18-\tfrac25\delta,\qquad X_2=48+\tfrac35\delta.
\]
Both quantities remain nonnegative for
\(-80\leq\delta\leq45\), the interval over which the objective is
\(3780+16\delta\). At the endpoints the basis structure changes;
outside the interval the problem must be solved again. This derivation is an
academic calculation for the example: QSB currently displays neither shadow
prices nor sensitivity ranges in the specific LP output or its native view.

\section{Method, backends, and limitations}
The \texttt{native} backend uses two-phase simplex: it first seeks a feasible
solution and then improves the objective. It is deterministic and educational,
using floating-point arithmetic. A mathematical optimum does not imply
symbolic precision. Widely differing coefficient scales can impair
calculations; use consistent units and check residuals.

\texttt{validate} runs the checks without computing an optimum.
\texttt{external} requests HiGHS, installed independently and discovered
through \texttt{QSB\_HIGHS\_PATH} or \texttt{PATH}; if it is unavailable, the
request fails explicitly. There is no silent fallback to the native backend.
The educational engine does not provide the preprocessing tools or scale
guarantees of an industrial engine. To exchange the model with another
optimizer, you can generate MPS:
\begin{lstlisting}
swift run qsb export-mps manual/examples/lp_sample.json > /tmp/qsb-lp.mps
\end{lstlisting}
MPS preserves names through a deterministic adaptation and expresses bounds
and domains; JSON remains QSB's normalized representation.

Distinguish three situations: an invalid model contains inconsistent data;
an infeasible model has no decision satisfying all constraints; an unbounded
model admits unlimited improvements in the objective. Adding an arbitrary
bound to make an error disappear changes the model's meaning. Review the
units, signs, and limits that the case actually imposes.

\section{Exercises}
\begin{enumerate}
 \item Increase the capacity of resource 1 to 190. Before running, predict
 \(X_1=14\), \(X_2=54\), and objective 3940 using the sensitivity analysis above;
 confirm the constraints with QSB.
 \item Reduce the capacity of resource 2 to 120 and solve. Explain why a
 marginal valuation should not be extrapolated indefinitely.
 \item Add \(X_1\geq60\). Identify the contradiction with resource 2 and
 compare a structurally valid model with an infeasible problem.
 \item Change the variables to integers and study the next chapter.
 Why does the optimum remain unchanged in the original example?
\end{enumerate}

\chapter{Integer and mixed-integer programming}
\label{chap:entera}

\section{Indivisible decisions}
Integer programming requires some decisions to take discrete values:
selected projects, shifts, batches, or open facilities. A mixed-integer linear
model retains a linear objective and linear constraints but combines
continuous and integer variables \cite{nemhauser1988integer} and \cite{wolsey1998integer}.
Integrality is part of the model, not a presentation operation on the result.
Rounding a continuous optimum can violate constraints or lose optimality.

An integer variable represents indivisible quantities; a binary variable
represents a yes/no decision. In QSB the normalized types are \texttt{continuous},
\texttt{integer}, and \texttt{binary}. All share the coefficients, relations, and
bounds of the linear model. In a selection model, a binary \(y_j\) can activate
a continuous quantity \(x_j\) through \(0\leq x_j\leq M_jy_j\).
The constant \(M_j\) must be a justifiable bound: an excessively large number
weakens the relaxation and can impair numerical stability.

\section{Example: minimum coverage with integer activities}
Two activities provide coverage for two requirements. Each unit of the
first costs 2.5 monetary units, and each unit of the second costs 2.
Their contributions are \((6,3)\) and \((3,5)\); at least 200 units of the
first requirement and 180 of the second are needed. Activities can only be
purchased in whole units:
\[
\begin{aligned}
 \min\quad &2.5X_1+2X_2\\
 \text{subject to}\quad &6X_1+3X_2\geq200,\\
 &3X_1+5X_2\geq180,\\
 &X_1,X_2\in\Z_+.
\end{aligned}
\]
Cost is additive and contributions remain constant: discounts and changes in
productivity have not been modeled. The complete case is in
\path{manual/examples/ilp_sample.json}; both variables have
\texttt{variableTypes} set to \texttt{integer} and a lower bound of zero.

\subsection{CLI and macOS execution}
From the repository root:
\begin{lstlisting}
swift run qsb solve-json-ilp manual/examples/ilp_sample.json --backend native
\end{lstlisting}
Use the specific ILP command to respect the integer domains; solving only
the relaxation does not answer the indivisible decision problem.

On macOS, use the same \texttt{Linear / Integer Programming} editor as in
Chapter~\ref{chap:pl}. Open the JSON or enter the two variables with
\texttt{Minimize}, coefficients 2.5 and 2, type \texttt{Integer}, a lower bound
of zero, and no explicit upper bound. Enter the two constraints with
relation \(\geq\), validate, and press \texttt{Run ILP}. The button label changes
according to the variable types. The solution displays quantities and
constraint evaluations; retain the exported JSON to review the original
domains.

\subsection{Optimum and verification}
QSB returns \(X_1=22\), \(X_2=23\), and cost 101. Feasibility is verified by
\[
 6(22)+3(23)=201\geq200,\qquad3(22)+5(23)=181\geq180.
\]
Coverage surpluses are one unit for each requirement. The cost is
\(2.5(22)+2(23)=101\), and both quantities are integers.

To understand the difference from the relaxation, temporarily allow
continuous variables. The intersection of the two boundaries is
\[
 X_1=\frac{460}{21}\approx21.904762,\qquad
 X_2=\frac{160}{7}\approx22.857143,
\]
with cost \(2110/21\approx100.476190\). In minimization, relaxing integrality
can only reduce the optimal cost or leave it unchanged; it provides a lower
bound for the integer problem. The absolute difference is
\(11/21\approx0.523810\). Rounding up happens to produce the optimal integer
solution here, but this is a coincidence of the case, not a general rule.

There is also a short verification of the integer cost. For any feasible
solution, the cost is at least \(2110/21\) and lies on the grid of multiples
of 0.5. Ruling out cost 100.5 is therefore sufficient to prove that 101 is the
minimum. If the cost were 100.5, then \(5X_1+4X_2=201\). Substitution into
the two coverage requirements simultaneously yields
\(X_1\geq197/9\) and \(X_1\leq285/13\); the narrow interval
\([21.888889,21.923077]\) contains no integer. Thus, cost 100.5 is
unattainable, and QSB's solution certifies the next possible value, 101.

\section{Branch and bound}
The native algorithm solves an LP relaxation at each node. If the solution
is integral, it becomes a candidate. If a variable takes a fractional value \(v\),
two problems are created, \(x_j\leq\lfloor v\rfloor\) and
\(x_j\geq\lceil v\rceil\), covering its integer possibilities. An infeasible
branch is discarded; a branch whose bound cannot improve the best known
solution is also discarded. Exploring all relevant branches establishes
optimality within numerical tolerances.

QSB uses a deterministic search and a default limit of 10000 nodes.
If this is exceeded, it reports a model unsupported by that search; this
must not be interpreted as an optimal solution or proof of infeasibility.
The educational backend lacks the cutting-plane and preprocessing techniques
of an industrial MILP engine. Running time can grow rapidly even for a small
file; providing valid physical bounds can improve the formulation but must
never exclude feasible decisions.

\texttt{validate} checks the model without solving it.
\texttt{external} uses the HiGHS executable discovered on the computer for LP/MIP
and requires an independent installation. The selected backend does not change
the meaning of the variables; if unavailable, the request fails.
The \texttt{export-mps} command accepts this JSON and expresses its integer
domains, bounds, and linear rows. This MILP capability does not imply that
HiGHS supports mixed-integer quadratic programming; these are different classes.

\section{Interpretation and common mistakes}
Compare costs only between feasible solutions of the same model. The
relaxation bound measures the idealized cost when fractions are allowed,
but does not represent an implementable alternative. In binary models, check
that each value represents the intended action; in mixed models, distinguish
continuous quantities from activation decisions. A fractional solution obtained
by changing the types to continuous answers a different problem.

Do not multiply a binary variable by a continuous variable within an LP row:
that is a nonlinear expression. Activation rules require auxiliary linear
constraints. Likewise, do not convert a continuous variable to integer simply
to make the screen display round numbers: the domain reflects the decision
unit.

\section{Exercises}
\begin{enumerate}
 \item Duplicate the example JSON and change both types to
 \texttt{continuous}. Solve with \texttt{solve-json} and contrast the lower
 bound calculated above with the integer result.
 \item Impose \(X_1\leq21\). Calculate the minimum \(X_2\) satisfying both
 coverage requirements, and check the new cost with QSB.
 \item Formulate a selection problem with three projects, costs 4, 6, and 7,
 benefits 5, 8, and 10, budget 10, and binary variables. Determine the optimal
 combination by enumeration, then verify it.
 \item Explain why rounding down both quantities in the example relaxation
 violates the requirements. What would change if the constraints were maximum
 capacities?
\end{enumerate}

\chapter{Goal, quadratic, and nonlinear programming}
\label{chap:optimizacion-avanzada}

This chapter brings together three distinct extensions of linear programming.
Goals order multiple objectives; quadratic programming adds curvature to the
objective while retaining linear constraints; nonlinear programming allows
general expressions. Choosing a family changes both the formulation and the
guarantee the solver can offer. Comparing the number displayed as the objective
is insufficient: priorities, feasibility, and exactness must also be reviewed.

\section{Goal programming}
\subsection{Priorities and deviations}
An organization may first seek to meet a goal and then optimize another
quantity without sacrificing what has already been achieved. In lexicographic
programming, \(g_1(x)\) is optimized, its optimal value is fixed, and
\(g_2(x)\) is optimized within that set; the process continues in the declared
order. Objectives are not automatically added with large weights. The order
expresses a policy decision: an improvement in a later priority never
compensates for deterioration in an earlier one.

For a goal \(a^\top x=t\), the deviation variables satisfy
\[
 a^\top x+n-p=t,\qquad n,p\geq0.
\]
\(n\) measures shortfall and \(p\) surplus. Minimizing \(n+p\) penalizes both
directions; minimizing only \(n\) tolerates surplus. Coefficients and units
must express the intended penalty. QSB does not infer these equations:
the user declares variables, objectives, and constraints.

\subsection{Example: protecting a production commitment}
There is capacity for ten units shared between products \(A\) and \(B\).
The first aim is to bring \(A\) close to eight units, followed by maximizing \(B\):
\[
\begin{aligned}
 &\text{Priority 1: }\min(n_1+p_1),\\
 &\text{Priority 2: }\max B,\\
 &A+B\leq10,\qquad A+n_1-p_1=8,\\
 &A,B,n_1,p_1\geq0.
\end{aligned}
\]
The priority does not itself require exact fulfillment if the goal is
impossible; it minimizes deviation subject to the hard constraints.
In this example, zero deviation is attainable. Once it is fixed,
\(A=8\), so the subsequent maximum is \(B=2\).

The file \path{manual/examples/opt_goal.json} contains all variables,
the two rows, and the ordered \texttt{goals} list; each priority has a name,
\texttt{sense}, and \texttt{coefficients}. Bounds and types are declared in
arrays aligned with \texttt{variableNames}. Run:
\begin{lstlisting}
swift run qsb validate-goal-json manual/examples/opt_goal.json
swift run qsb solve-goal-json manual/examples/opt_goal.json --backend native
\end{lstlisting}
The output document retains the model, \texttt{backend} metadata, and
solution.

\texttt{goalOutcomes} displays priority, sense, and value.
\begin{table}[htbp]
\centering
\caption{Lexicographic result for the production case}
\begin{tabular}{lr}
\toprule
Quantity & Value \\
\midrule
Priority 1: total deviation & 0 \\
Priority 2: production of \(B\) & 2 \\
\(A\) & 8 \\
\(n_1\), \(p_1\) & 0, 0 \\
\bottomrule
\end{tabular}
\end{table}
Capacity is active and the goal is met. Reversing the order makes the first
objective yield \(B=10\); with that value fixed, \(A=0\), \(n_1=8\), and
deviation is eight. This contrast invites discussion of who determines
precedence and which commitment it protects.

\subsection{macOS workflow and solver character}
Open the normalized file in \texttt{Goal Programming}. In
\texttt{Definition}, review variables and bounds, the priority list, and
linear constraints. The move-up and move-down controls change a priority's
position. Enter the deviation equation as an equality and verify the
negative sign of \(p_1\). Press \texttt{Validate} and \texttt{Run}; the native
priority view shows the ordered objectives, final quantities, and backend
information. JSON provides the complete exportable record. Tables support
copying and pasting; coefficient columns must match the current variable
order. Large tables use row selection and coefficient groups.

The native backend solves a sequence of LP or integer problems, adding an
equality to preserve each previous optimum. The continuous case is labeled
\texttt{exact}; this term characterizes the method, subject to floating-point
tolerances. Integer goals use an educational search for small problems,
identified as \texttt{fixtureScale}, and depend on branch-and-bound limitations.
The special relaxation of integer deviation pairs applies only when integral
equations justify integral residuals; it does not remove integrality from
the decisions.

\texttt{validate} does not solve. \texttt{external} executes each priority
with HiGHS if installed; QSB preserves the lexicographic order and fixing.
A missing executable causes an explicit error. An infeasible later priority,
an unbounded objective, or exhausted limits require reviewing the diagnostic
rather than interpreting partial results as the final plan.

\section{Quadratic programming}
\subsection{Objective convention and assumptions}
QSB represents the objective as
\[
 f(x)=x^\top Qx+c^\top x,
\]
with linear constraints, bounds, and variable types. JSON uses
\texttt{quadraticMatrix} for \(Q\) and \texttt{linearCoefficients} for \(c\).
This convention has no factor of \(1/2\). For a symmetric matrix, the
coefficient of the product \(x_ix_j\), \(i\ne j\), is \(2Q_{ij}\): a term
\(4XY\) is represented by \(Q_{12}=Q_{21}=2\). Do not place 4 in both
cells, since that would double the effect.

A symmetric positive-definite matrix produces a strictly convex objective;
its minimum over a convex region, if it exists, is unique. A negative-definite
matrix produces the equivalent strictly concave maximization case.
Stationarity and feasibility conditions then identify the global optimum.
An indefinite matrix changes the problem and is not covered by that guarantee.

\subsection{Example: adjustment with limited capacity}
Two quantities should be brought close to 2 and 3, respectively, with total
capacity four. Removing the constant 13 from the squared distance gives
\[
\begin{aligned}
 \min\quad & X^2+Y^2-4X-6Y\\
 \text{subject to}\quad &X+Y\leq4,\qquad X,Y\geq0.
\end{aligned}
\]
The matrix is the identity and the linear vector is \((-4,-6)\). Removing
the constant preserves optimal quantities but changes the reported objective
value: the original squared distance is \(f(X,Y)+13\).

The model is in \path{manual/examples/opt_qp.json}. Run:
\begin{lstlisting}
swift run qsb validate-qp-json manual/examples/opt_qp.json
swift run qsb solve-qp-json manual/examples/opt_qp.json --backend native
\end{lstlisting}
The optimum is \(X=1.5\), \(Y=2.5\), \(f=-12.5\), with
\texttt{Capacity} in \texttt{activeConstraints}. The original squared distance
is 0.5. A negative objective is perfectly valid: it results from the omitted
constant and does not represent a negative distance.

Without the capacity constraint the solution would be \((2,3)\), which is
infeasible. With multiplier \(\lambda\geq0\), the stationarity equations are
\[
 2X-4+\lambda=0,\qquad2Y-6+\lambda=0.
\]
The active capacity requires \(X+Y=4\), giving \(\lambda=1\), \(X=1.5\), and
\(Y=2.5\). The lower bounds are inactive. Strict convexity and the conditions
above certify the global minimum of the continuous model.
QSB displays active constraints, but this multiplier derivation is an
independent academic check.

\subsection{macOS workflow and limitations}
Open the JSON in \texttt{Quadratic Programming}. \texttt{Definition} separates
variables and bounds, the quadratic matrix, and linear constraints. Enter
the linear coefficients in the variable rows and the identity in the matrix;
confirm the names on both axes. Add \texttt{Capacity} with coefficients
\((1,1)\), relation \(\leq\), and right-hand side 4. Validate and press
\texttt{Run}. The native solution displays the objective, variables, bounds,
active constraints, and exactness; JSON allows inspection of the document.

\texttt{Copy Table} and \texttt{Paste Table} transfer the complete matrix.
Reordering known variables preserves both axes by name; pasted headers must
match the current order. For more than twelve variables, the editor displays
one row and groups of columns, although table exchange remains complete.
This editor capability does not imply support for large problems in the
native solver.

The continuous native solver enumerates active sets and solves KKT systems
for strictly convex minimization or strictly concave maximization. Its label
is \texttt{exact}, with numerical calculations. The search rejects more than
twenty candidate inequalities, counting bounds that it converts into
inequalities; the cost of set enumeration grows rapidly. The integer native
solver supports only diagonal concave maximization with nonnegative variables,
up to four variables, and enumerates points within derived finite bounds.
It is identified as \texttt{fixtureScale}; it is not a general MIQP solver.

The external HiGHS backend supports translation of the supported strictly
convex/concave continuous cases. QSB adapts the convention using a Hessian
\(H=2Q\), because the external solver uses \(\tfrac12 x^\top Hx\).
The adaptation also preserves sense, bounds, and constraints, and verifies
the objective and feasibility when reading the solution.
HiGHS does not solve MIQP in this integration; quadratic integer variables
require the native scope described above or a future engine. The
\texttt{validate} mode does not compute an optimum.

\section{Nonlinear programming}
\subsection{Expressions and domain}
Nonlinear programming (NLP) allows an objective \(f(x)\) and constraints
\(g_i(x)\ \{\leq,=,\geq\}\ b_i\). QSB accepts textual expressions,
ordered variables, and nullable bounds. The evaluator recognizes addition,
subtraction, multiplication, division, powers, parentheses, implicit
multiplication, and \texttt{exp}, \texttt{log}, \texttt{sqrt}, \texttt{sin},
\texttt{cos}. Use explicit multiplication in teaching materials to aid reading.
Names are recognized without case sensitivity; consistent spelling is advisable.

Valid syntax does not guarantee a valid domain at every point:
\texttt{log(X)} requires \(X>0\), \texttt{sqrt(X)} requires \(X\geq0\), and
division requires a nonzero denominator. Bounds must correspond to the case
and protect the domain where possible. Sine and cosine angles are expressed
in radians. JSON stores \texttt{objectiveExpression} and, for each constraint,
\texttt{expression}, \texttt{relation}, and \texttt{rhs}.

\subsection{Example: constrained distance}
The file \path{manual/examples/opt_nlp.json} specifies
\[
 \min\ (X-3)^2+(Y-2)^2,\qquad X+Y\leq4,\qquad0\leq X,Y\leq4.
\]
The objective is the squared distance to point \((3,2)\), which violates
capacity. Its projection onto \(X+Y=4\) is \((2.5,1.5)\), and the mathematical
value is 0.5. Although this case could be formulated as QP, it serves to
compare expression input and an approximate method with a known analytical
solution.
\begin{lstlisting}
swift run qsb validate-nlp-json manual/examples/opt_nlp.json
swift run qsb solve-nlp-json manual/examples/opt_nlp.json --backend native
\end{lstlisting}
The verified execution returns approximately
\(X=2.500000002\), \(Y=1.499999999\), objective 0.4999999995, and maximum
violation \(5.0\times10^{-10}\). The evaluation of \texttt{Capacity} is
approximately 4.0000000005. The \texttt{iterations} field records 361 cumulative
iterations of the selected search; it does not count all evaluations of
searches from the other initial points. This value describes this execution,
not a universal limit or a measure comparable with simplex iterations.

The small violation does not authorize physically exceeding capacity. It is
a numerical search residual; for real implementation, check the values
against the tolerance and units of the case. The output contains
\texttt{constraintEvaluations}, \texttt{maximumViolation},
\texttt{iterations}, the objective, and variable values, as well as the model
and metadata. Reporting the objective without the residual would hide
relevant information.

\subsection{macOS use and approximate guarantee}
Open the JSON in \texttt{Nonlinear Programming}. The definition allows editing
the sense, objective expression, bounded variables, and constraints with an
expression, relation, and right-hand side. Enter
\texttt{(X-3)\^{}2+(Y-2)\^{}2}, bounds 0 and 4, and \texttt{X+Y} with
\(\leq4\). Use \texttt{Validate} before \texttt{Run}. An invalid expression is
retained for correction. Variables and constraints support spreadsheet
exchange; each valid paste replaces a complete table.
The solution displays the candidate, bounds, evaluations, and maximum
violation, together with its approximate character and exportable JSON.

The native backend uses automatic differentiation, searches from multiple
points, and progressive constraint penalties. It is deterministic, is limited
to four variables, and is identified as \texttt{approximate}. It also includes
a low-dimensional search on equality manifolds for small cases.
These mechanisms do not prove global optimality for nonconvex functions.
A feasible candidate may be local, and a failed search does not in itself
constitute a mathematical proof of infeasibility.

Strict inequalities from imports are normalized to non-strict boundaries
and generate a structured warning. The
\texttt{normalizedStrictInequalities} field records that transformation;
it must be reviewed if the boundary has a physical meaning excluded by the
case. \texttt{validate} diagnoses without solving. \texttt{external} is
explicitly unavailable for NLP; select the native backend for supported
educational problems.

\section{How to compare and review results}
Retain each input together with its output and the backend used. For goals,
review the order and each fixed value. For QP, check the matrix convention
and supported curvature. For NLP, review domains, bounds, evaluations, and
maximum violation before describing a solution as implementable. File
validation and search convergence provide different kinds of evidence.

All three families have a native definition view and a JSON view; unfinished
fields remain editable when validation fails. In native macOS runs, solving
takes place in the background and duplicate runs are prevented while one is
in progress. Changing the model can make that calculation's result stale.
There is currently no cancellation control for these workers. Accessible
labels on controls and results do not replace a complete VoiceOver and
keyboard navigation audit, which remains pending for these editors.

\section{Exercises}
\begin{enumerate}
 \item Reverse the priorities in the goal example. Verify \(B=10\),
 \(A=0\), and total deviation 8. Explain which commitment has changed.
 \item Reduce the capacity in the goal example to 6. Calculate the minimum
 deviation of the first priority and justify why the second cannot recover
 production of \(B\) without worsening the first.
 \item In the QP, increase capacity to 5. Predict \((X,Y)=(2,3)\),
 objective \(-13\), and zero original distance. Check against QSB.
 \item Introduce a term \(4XY\) in a test objective. Write the correct
 symmetric matrix and examine its curvature before choosing a solver.
 \item Reformulate the example NLP objective in QP form, retaining
 constraints and bounds. Compare quantities and residuals; explain the
 difference between its objective with a constant and the quadratic
 objective without it.
 \item Propose a one-variable function with two local minima. Explain what
 additional evidence you would need to attribute global optimality to a
 candidate obtained by the native NLP solver.
\end{enumerate}

\chapter{Network models}
\label{chap:redes}

\section{From a situation to a graph}
A network represents places, states, or people as nodes; arcs represent
possible connections. The meaning of the number on an arc depends on the
model: cost or distance in a route, capacity in a maximum-flow problem.
Confusing these quantities changes the problem \cite{ahuja1993network}.
The examples in this chapter are original, and their files are in
\texttt{manual/examples/}. They all use the JSON envelope with \texttt{kind}
and \texttt{model}; WinQSB files are not required.
Run the commands from the repository root.

\section{Shortest path and minimum spanning tree}
SPP seeks a directed path from the source to the destination with minimum cost.
Dijkstra requires nonnegative costs. In QSB, the first node is the source and
the last is the destination. An absent arc is not equivalent to a zero-cost arc.
The file \path{plan_spp.json} has nodes A, B, C, D and the following arcs:
\begin{center}
\begin{tabular}{lrrrrr}\toprule
Arc & A--B & A--C & B--C & B--D & C--D \\
Cost & 2 & 5 & 1 & 6 & 2 \\
\bottomrule\end{tabular}
\end{center}
\begin{lstlisting}[language=bash]
swift run qsb validate-network-json manual/examples/plan_spp.json
swift run qsb solve-network-json manual/examples/plan_spp.json --backend native
swift run qsb solve-network-json manual/examples/plan_mst.json --backend native
\end{lstlisting}
\begin{figure}[htbp]
\centering
\begin{tikzpicture}
\node[qsb node] (A) at (0,0) {A};
\node[qsb node] (B) at (3,1.2) {B};
\node[qsb node] (C) at (3,-1.2) {C};
\node[qsb node] (D) at (6,0) {D};
\draw[qsb selected] (A) -- node[qsb label,above] {2} (B);
\draw[qsb alternative] (A) -- node[qsb label,below] {5} (C);
\draw[qsb selected] (B) -- node[qsb label,right] {1} (C);
\draw[qsb alternative] (B) -- node[qsb label,above] {6} (D);
\draw[qsb selected] (C) -- node[qsb label,below] {2} (D);
\end{tikzpicture}

\caption{Shortest-path network: labels are costs. Thick arrows trace the optimal route; dashed arrows show alternatives.}
\label{fig:shortest-path}
\end{figure}

SPP returns A $\to$ B $\to$ C $\to$ D with cost $2+1+2=5$.
MST connects all nodes without cycles at the lowest total cost; it does not
seek a path between endpoints. The example uses the same connections as
undirected edges, and Kruskal chooses B--C, A--B, C--D, with cost 5.
QSB normalizes asymmetric MST connections to the lower cost for each pair;
a disconnected network has no tree spanning all nodes.

\section{Maximum flow}
With capacities $u_{ij}$, flow must satisfy
$0\leq x_{ij}\leq u_{ij}$ and conservation at intermediate nodes. The
objective maximizes net flow leaving the source. It does not minimize cost.
In \path{plan_maxflow.json}, S--A has capacity 3, S--B 2,
A--B 1, A--T 2, and B--T 3. Although the JSON field is named \texttt{cost},
it represents capacity in MFP.
\begin{lstlisting}[language=bash]
swift run qsb solve-network-json manual/examples/plan_maxflow.json --backend native
\end{lstlisting}
\begin{figure}[htbp]
\centering
\begin{tikzpicture}
\node[qsb node] (S) at (0,0) {S};
\node[qsb node] (A) at (3,1.2) {A};
\node[qsb node] (B) at (3,-1.2) {B};
\node[qsb node] (T) at (6,0) {T};
\draw[qsb arc] (S) -- node[qsb label,above] {3/3} (A);
\draw[qsb arc] (S) -- node[qsb label,below] {2/2} (B);
\draw[qsb arc] (A) -- node[qsb label,right] {1/1} (B);
\draw[qsb arc] (A) -- node[qsb label,above] {2/2} (T);
\draw[qsb arc] (B) -- node[qsb label,below] {3/3} (T);
\draw[densely dotted,line width=1pt] (0.75,-1.7) -- (0.75,1.7);
\end{tikzpicture}

\caption{Maximum flow: each label gives flow/capacity. The dotted line separates the source from the remaining nodes and marks a saturated cut.}
\label{fig:maximum-flow}
\end{figure}

Edmonds--Karp obtains flow 5: 3 and 2 units leave S, and 2 and 3 reach T.
The cut separating S from the other nodes has capacity 5 and proves optimality.
Again, the first and last nodes determine the endpoints.

\section{Traveling salesperson}
TSP seeks a directed cycle that visits every node exactly once and returns
to the starting node. Connecting the nodes is insufficient: there must be
a single tour. The example \path{plan_tsp.json} contains three nodes, with
$c_{AB}=2$, $c_{BC}=3$, $c_{CA}=4$ and reverse costs 4, 5, and 6.
\begin{lstlisting}[language=bash]
swift run qsb solve-network-json manual/examples/plan_tsp.json --backend native
\end{lstlisting}
The tour A $\to$ B $\to$ C $\to$ A costs 9; the reverse tour costs 15.
Held--Karp is exact for the accepted model but exponential; the solver
rejects more than 20 nodes. A model with no feasible Hamiltonian cycle is
not repaired by adding connections. The \texttt{fixtureScale} label expresses
this scale limitation, not a guarantee of fast execution near the maximum.

\section{Assignment and transportation}
Assignment minimizes $\sum_{ij}c_{ij}x_{ij}$, assigning each worker to one
task and using each task at most once. The Hungarian method accepts
rectangular matrices with at least as many tasks as workers.
In \path{plan_assignment.json}, the rows Ana and Luis have costs
$(2,7,4)$ and $(6,1,3)$ for X, Y, and Z.
\begin{lstlisting}[language=bash]
swift run qsb solve-network-json manual/examples/plan_assignment.json --backend native
swift run qsb solve-network-json manual/examples/plan_transport.json --backend native
\end{lstlisting}
Ana--X and Luis--Y cost 3; Z remains unassigned. The model does not include
teams, multiple tasks per person, or additional constraints.

Transportation allows continuous quantities and requires balance:
\[
 \min\sum_{ij}c_{ij}x_{ij},\qquad
 \sum_jx_{ij}=s_i,\quad \sum_ix_{ij}=d_j,\quad x_{ij}\geq0.
\]
The example has supply $(4,6)$, demand $(5,5)$, and cost matrix
$\left(\begin{smallmatrix}1&4\\3&2\end{smallmatrix}\right)$.
QSB ships 4 from O1 to D1, 1 from O2 to D1, and 5 from O2 to D2;
the cost is $4+3+10=17$. Check row and column sums separately.
An imbalance is a validation error in TP. This family does not automatically
create a dummy source or destination.

\section{Minimum-cost flow}
CNF combines transshipment and costs: at each node,
$\sum_jx_{ij}-\sum_jx_{ji}=s_i-d_i$, and it minimizes
$\sum_{ij}c_{ij}x_{ij}$ with nonnegative continuous quantities.
The current model has no arc capacities or lower bounds.
The validator requires finite nonnegative costs in CNF and the graph
variants; assignment permits finite negative costs. Do not model a subsidy
as a negative cost on a CNF arc, because it will be rejected.
\path{plan_netflow.json} supplies 4 at O, requires 4 at D, and allows
O--M at cost 1, M--D at cost 2, and O--D at cost 5.
\begin{lstlisting}[language=bash]
swift run qsb solve-network-json manual/examples/plan_netflow.json --backend native
\end{lstlisting}
The 4 units pass through M and cost 12; direct shipment would cost 20.
CNF preserves original supply and demand and, when they differ, emits a
warning and adds a zero-cost dummy adjustment during solving.
Always review \texttt{balanceAdjustments}: a solution with an adjustment
does not prove that original supply can meet actual demand.

\section{macOS use and backend selection}
Open the example by importing normalized JSON. In Definition,
SPP, MST, MFP, and TSP have a node and arc editor;
AP and TP have cost matrices, with supply and demand in TP; CNF has nodes,
arcs, and balances. Changes are validated before Run. The endpoint order
for SPP/MFP maps to the JSON order; visual positions are not exported.
The solution Diagram highlights the path, tree, tour, or connections used.
\texttt{All connections} also shows the remaining connections. In bipartite
networks, the groups represent workers/tasks or sources/destinations.
Use the detail panel to read exact values and JSON to audit the result;
the thickness of a connection does not replace its numerical quantity.

Native offers deterministic algorithms for all seven variants. Validate
checks the model without solving it. External is available for TP and CNF
through HiGHS; there is no external backend covering all network variants.
Solutions preserve the model, algorithm, assumptions, and exactness. The
practical size of native LPs depends on the formulation; for large
transportation networks, consider HiGHS and check its availability.

\section{Exercises}
\begin{enumerate}
\item Change A--C from 5 to 1 in SPP. Verify the new path and explain why
      the minimum spanning tree may respond differently.
\item Reduce S--A to 1 in MFP. Calculate a cut before solving.
\item Change D2 demand to 6 in TP. Observe validation; then
      deliberately represent one additional unit of supply.
\item In CNF, increase demand to 5. Interpret the dummy adjustment and
      explain which business data are missing before the plan can be accepted.
\end{enumerate}

\chapter{Decisions, games, and dynamic programming}
\label{chap:decisiones}

\section{Choosing under uncertainty}
A decision combines controllable alternatives with states of the environment
that the analyst does not control. A payoff table organizes the payoff $a_{ij}$
from alternative $i$ in state $j$. When prior probabilities $p_j$ are known,
the expected-value criterion selects
\[
 i^*\in\arg\max_i\sum_j p_j a_{ij}.
\]
This criterion assumes comparable payoffs and risk neutrality. A solution
maximizing average profit may involve substantial losses in a particular state.
QSB does not automatically transform money into utility or incorporate the
decision maker's risk aversion.

\subsection{Payoff tables and information}
The file \path{dec_payoff.json} considers a launch with demand Alta (High,
probability $0.4$) or Baja (Low, $0.6$). The input's Spanish identifiers are
retained here so that they match the files and output: Lanzar means Launch,
Esperar means Wait, Positivo means Positive, and Negativo means Negative.
The payoffs are:
\begin{center}
\begin{tabular}{lrr}
\toprule
Decision & Alta & Baja\\
\midrule
Lanzar & 100 & $-40$\\
Esperar & 20 & 10\\
\bottomrule
\end{tabular}
\end{center}
Without a market study, Lanzar yields $0.4(100)+0.6(-40)=16$, compared with
14 for Esperar. With perfect information, Lanzar would be chosen in Alta and
Esperar in Baja: $0.4(100)+0.6(10)=46$. The expected value of perfect
information (EVPI) is $46-16=30$, an upper bound on the price of information
under this model.

The study produces a Positivo or Negativo indicator. The
\texttt{indicatorLikelihoods} matrix uses indicator rows and state columns:
$P(+\mid Alta)=0.8$ and $P(+\mid Baja)=0.2$. Each column represents a
distribution over possible indicators. These are not posterior probabilities.
The policy with the study is Lanzar if the result is Positivo and Esperar if
it is Negativo. Its expected value is $33.6$; hence EVSI (expected value of
sample information) is $33.6-16=17.6$. This amount does not deduct the cost
of the study. If the study cost 20, purchasing it would not be worthwhile
under the expected-value criterion.

\begin{lstlisting}
swift run qsb validate-decision-json manual/examples/dec_payoff.json
swift run qsb solve-decision-json manual/examples/dec_payoff.json --backend native
\end{lstlisting}
The output contains \texttt{priorExpectedValues}, \texttt{indicatorAnalyses},
\texttt{bestPriorDecision}, EVSI, and EVPI. Keys retain their English names.
The solver requires a probability greater than $10^{-12}$ for every indicator;
otherwise, solution fails rather than inventing a posterior.

\subsection{Bayesian updating}
Bayes' rule converts prior probability and likelihood into posterior probability:
\[
 P(s_j\mid e_k)=\frac{P(e_k\mid s_j)p_j}{\sum_\ell P(e_k\mid s_\ell)p_\ell}.
\]
The file \path{dec_bayes.json} separates this operation from any payoff table.
It uses the same states and study. QSB obtains $P(+)=0.44$, $P(-)=0.56$,
and the following posteriors:
\begin{center}
\begin{tabular}{lrr}
\toprule
Evidence & Alta & Baja\\
\midrule
Positivo & $0.727273$ & $0.272727$\\
Negativo & $0.142857$ & $0.857143$\\
\bottomrule
\end{tabular}
\end{center}
\begin{lstlisting}
swift run qsb validate-decision-json manual/examples/dec_bayes.json
swift run qsb solve-decision-json manual/examples/dec_bayes.json
\end{lstlisting}
\texttt{likelihoods} again uses outcome rows and state columns. Changing the
order of names requires changing the corresponding rows or columns as well.
Probabilities must be finite, nonnegative, and consistent with the distributions
represented. Bayes updates beliefs; by itself it does not say which action to take.

\subsection{Trees and backward induction}
A tree makes the temporal order of decisions and observations explicit.
\texttt{rootID} identifies the root, each node has a unique \texttt{id}, and
its \texttt{childIDs} identify successors. The types are \texttt{decision},
\texttt{chance}, and \texttt{terminal}. Chance-branch probabilities are stored
in the child nodes, rather than in a list on the parent.

Backward induction assigns a terminal its payoff. A decision node adds its
local payoff to the maximum value of its children; a chance node adds its
local payoff to the expected value of its children. Local payoffs accumulate
along the path: do not repeat a cost in the leaves if it was already entered
in an ancestor.

In \path{dec_tree.json}, Invertir (Invest) costs 5 and yields 40 with
probability $0.6$ or $-10$ with probability $0.4$. Its value is
$-5+0.6(40)+0.4(-10)=15$; Conservar (Retain) yields 12. The root value is
15, and the policy selects Invertir.
\begin{lstlisting}
swift run qsb validate-decision-json manual/examples/dec_tree.json
swift run qsb solve-decision-json manual/examples/dec_tree.json
\end{lstlisting}
\begin{figure}[H]
\centering
\begin{tikzpicture}
\node[rectangle,draw,minimum size=6mm,label=above:{\small Elegir}] (R) at (0,0) {};
\node[qsb node,label=above:{\small Invertir},label=below:{$-5$}] (I) at (4,1) {};
\node[qsb terminal] (K) at (4,-1.5) {Conservar\\12};
\node[qsb terminal] (E) at (8,2) {Exito\\40};
\node[qsb terminal] (F) at (8,-0.2) {Fracaso\\$-10$};
\draw[qsb selected] (R) -- (I);
\draw[qsb alternative] (R) -- (K);
\draw[qsb arc] (I) -- node[qsb label,above] {$0.6$} (E);
\draw[qsb arc] (I) -- node[qsb label,below] {$0.4$} (F);
\end{tikzpicture}

\caption{Decision tree: the square represents a choice and the circle represents chance. The thick branch is the selected policy. Leaves show payoffs; the local investment cost is added once. Names retain the JSON identifiers.}
\label{fig:decision-tree}
\end{figure}

Review \texttt{nodeValues} for each node and \texttt{policy} for the decisions.
The method maximizes payoffs: to minimize costs, you can use their negatives,
explaining the convention. QSB warns about and normalizes chance-child
probabilities whose positive sum differs from one, for compatibility with
rounded inputs. This repair does not replace a review of the data. Avoid cycles,
nonexistent identifiers, and inconsistent probabilities.

\section{Zero-sum games}
In a zero-sum game, the row player's payoff is the column player's loss. The
first maximizes and the second minimizes. Mixed strategies are probability
distributions over actions; they do not represent fractional quantities of an
action in a single play.

For a matrix $A$, QSB solves equivalent linear programs. The row player seeks
$\max v$ subject to $A^Tp\geq v\mathbf{1}$, $\sum_i p_i=1$, and $p_i\geq0$.
The column player solves the dual minimization problem. This assumes simultaneous
choices and known payoffs. \path{dec_game.json} uses
\[
 A=\begin{pmatrix}1&-1\\-1&1\end{pmatrix}.
\]
Both players assign probability $0.5$ to A and B, and the game's value is
zero. The numerical output may show $4.44\times10^{-16}$ because of rounding.
\begin{lstlisting}
swift run qsb validate-decision-json manual/examples/dec_game.json
swift run qsb solve-decision-json manual/examples/dec_game.json
\end{lstlisting}
Read \texttt{rowStrategy}, \texttt{columnStrategy}, and \texttt{value}. The
native method uses the shared linear solver. It does not interpret this solution
as an equilibrium of a general non-zero-sum game or a sequential game.

\section{Dynamic programming: state, stage, and action}
Dynamic programming divides a problem into stages. The state summarizes the
information needed to evaluate future decisions; an action changes that state
and produces a local cost or return. The principle of optimality allows an
optimal policy to be built from subproblems. A state omitting relevant
information breaks that principle.

QSB implements three discrete, deterministic models: bounded knapsack,
stage-by-stage acyclic routing, and production with inventory. Inputs use
\texttt{kind} and \texttt{model}. Outputs include a typed result,
\texttt{assumptions}, and \texttt{trace}. Each trace row describes a stage,
state, action, next state, and \emph{local} value; this value is not necessarily
the subproblem's optimal value function.

\subsection{Bounded knapsack}
Each item $i$ consumes $w_i$ units of capacity and contributes $r_i$ per unit;
at most $b_i$ units are available. The formulation is
\[
 \max\sum_i r_i x_i,\qquad \sum_i w_i x_i\leq C,
 \qquad x_i\in\{0,\ldots,b_i\}.
\]
A recurrence with remaining capacity $c$ is
$F_i(c)=\max_q\{q r_i+F_{i+1}(c-qw_i)\}$, over available quantities that fit.
Capacity and consumption are integers; returns may be finite real numbers.
Do not convert continuous physical measurements to integers without justifying
the discretization unit.

In \path{dec_knapsack.json}, $C=7$ and the items $(b_i,w_i,r_i)$ are
A $(2,3,5)$, B $(3,2,3)$, and C $(1,4,7)$. QSB selects one A and one C:
capacity 7 and return 12. The alternative of one A and two B also uses 7 but
returns 11. The selection encodes quantities; the trace also includes the
decision to take zero units of B.
\begin{lstlisting}
swift run qsb validate-dp-json manual/examples/dec_knapsack.json
swift run qsb solve-dp-json manual/examples/dec_knapsack.json
\end{lstlisting}
In macOS, Bounded Knapsack lets you edit capacity, availability, consumption,
and returns, and copy or paste the table. Summary provides metrics, bars, and
a selection table; the trace helps explain the policy.

\subsection{Stage-by-stage routing (Stagecoach)}
The \texttt{nodes} list declares a topological order: the first node is the
origin and the last the destination. Every arc must advance in that order and
have a finite, nonnegative cost. The method applies
\[
 V(i)=\min_{(i,j)}\{c_{ij}+V(j)\},\qquad V(\text{destination})=0.
\]
\path{dec_route.json} contains Origen--Norte (4), Origen--Sur (2),
Norte--Destino (3), and Sur--Destino (6). The identifiers mean Origin, North,
South, and Destination, respectively. The result is Origen--Norte--Destino,
with cost 7; choosing the cheapest initial arc would yield 8.
\begin{lstlisting}
swift run qsb validate-dp-json manual/examples/dec_route.json
swift run qsb solve-dp-json manual/examples/dec_route.json
\end{lstlisting}
A structurally valid graph may have no route to the destination; this is
detected during solution. This model does not accept cycles or backward arcs:
for other graphs, consult the networks chapter.

\subsection{Production and inventory}
The state at the beginning of a period is inventory $I_{t-1}$; the action is
integer production $P_t$. Deterministic demand $d_t$ is met with
$I_t=I_{t-1}+P_t-d_t\geq0$. Capacities bound $P_t$ and $I_t$. QSB assumes
zero initial and final inventory, with no backorders. It minimizes
\[
 \sum_t\left[K_t\mathbf{1}_{P_t>0}+c_tP_t+h_tI_t\right].
\]
Holding cost applies to the period's ending inventory, rather than its initial
inventory. The recurrence evaluates feasible production quantities and adds
the optimal future cost.

\path{dec_production.json} has demands of 2 and 3, production capacity 5,
storage capacity 3, setup cost 10, unit cost 2, and holding cost 1 in both
periods. QSB produces 5 initially, carries 3, and produces nothing in the
second period. The cost is $10+2(5)+1(3)=23$, compared with 30 if production
is set up and carried out in each period. The decisions detail initial
inventory, production, demand, ending inventory, and local cost.
\begin{lstlisting}
swift run qsb validate-dp-json manual/examples/dec_production.json
swift run qsb solve-dp-json manual/examples/dec_production.json
\end{lstlisting}
Field validation does not prove the existence of a feasible plan. Insufficient
capacities may cause solution failure. This model does not represent multiple
products, configurable initial inventory, shortages, or random demand.

\section{Use, metadata, and shared limits}
All four decision variants have native macOS editors. The tree editor allows
nodes and links to be edited; removing leaves or adding terminals maintains a
structure that is easier to review. Stagecoach and production/inventory can be
edited through the structured form or JSON; their results have native route or
per-period series views. Use Validation before Run and consult JSON when you
need to preserve every field.

The native decision backend declares \texttt{exact}: it computes the specified
criterion using numerical arithmetic within the model and its educational scale.
The three dynamic variants declare \texttt{fixtureScale}; their policies are
those obtained by the implemented recurrences, with discretization and memory
consumption growing with capacity, inventory, and the number of stages. There
is no promise of industrial-scale performance. \texttt{--backend validate}
returns validation rather than a solution; the external backend is unavailable
for these families. Save the model alongside the dynamic programming output:
its solution document preserves the title and assumptions, but not the complete
model.

\section{Exercises and checks}
\begin{enumerate}
 \item Change $P(Alta)$ from $0.4$ to $0.3$, adjusting Baja as well. First
 calculate the best decision without information, then solve again. Why can
 EVSI change even when the quality of the study remains unchanged?
 \item Calculate the maximum price of the example's study and compare it
 with EVPI. Explain which assumptions would prevent using that price in a
 real decision.
 \item In the tree, increase the cost of Invertir to 9. This branch's value
 becomes 11, and the root decision changes to Conservar, with value 12.
 \item In the knapsack, change capacity to 6. Verify a return of 10 with two
 units of A, and explain why the return/consumption ratio is insufficient to
 ensure optimality.
 \item In the route, make Sur--Destino cost 4. The route through Sur will
 cost 6. Then add a backward arc and examine the validation diagnostic.
 \item Reduce the production setup cost to 1 in both periods. Producing 2
 and 3 separately costs 12, compared with 14 when inventory is carried.
 Compare each inventory balance and the trace.
\end{enumerate}


\part{Operations planning}
\chapter{Inventory models}
\label{ch:inventarios}

\section{The decision being modeled}
An inventory policy combines how much to replenish and when to do so. The
minimum cost depends on time units, the demand pattern, and shortage costs;
entering average demand alone is not sufficient
\cite{harris1913many,silver1998inventory}. QSB distinguishes EOQ/EPQ,
quantity discounts, the newsvendor model, period-based lot sizing, and four
stochastic review policies. All use the \texttt{kind}/\texttt{model} JSON
envelope and the native or validation backend. There is no external inventory
backend.

\section{EOQ and gradual production}
With constant demand \(D\), fixed ordering cost \(S\), and holding cost
\(H\) per unit per period, with no shortages, the relevant cost is
\[
 C(Q)=DS/Q+HQ/2,\qquad Q^*=\sqrt{2DS/H}.
\]
Purchasing cost \(Dc\) is added to total cost but does not change the optimum
when unit price \(c\) is constant. Lead time \(\ell\) gives the reorder
point \(r=D\ell\). Rates, lead time, and holding cost must use the same
period. For gradual production at rate \(p>D\), average inventory is
\(Q(1-D/p)/2\), and
\[
 Q^*_{EPQ}=\sqrt{\frac{2DS}{H(1-D/p)}}.
\]
The calculation is continuous: purchasing boxes or whole units requires
comparing feasible quantities. QSB does not automatically impose packaging
constraints.

\subsection{Example: replenishing notebooks}
The file \path{manual/examples/prob_eoq.json} contains \(D=1200\)
units/year, \(S=30\), \(H=2\), \(c=5\), a lead time of \(0.01\)
years, and a current order quantity of \(100\). From the repository root:
\begin{lstlisting}
.build/debug/qsb validate-inventory-json manual/examples/prob_eoq.json
.build/debug/qsb solve-inventory-json manual/examples/prob_eoq.json
\end{lstlisting}
The results are \(Q^*=189.736660\), \(6.324555\) orders/year, a cycle of
\(0.158114\) years, and a reorder point of \(12\). Ordering and holding
each cost \(189.736660\); relevant cost is \(379.473319\), and total
cost is \(6379.473319\). Ordering \(100\) costs \(6460\): the model's
annual saving is \(80.526681\), rather than the entire purchasing cost.
The output document includes the comparison in \texttt{knownQuantity} and
\texttt{closedForm} exactness. In \path{prob_epq.json}, \(p=2400\)
increases the optimal lot to \(268.328157\), with relevant cost
\(268.328157\).

\section{All-units quantity discounts}
QSB evaluates each price tier and its feasible quantities, including purchasing,
setup, and holding costs. The discount applies to every unit in the order; it
is not a marginal block tariff. Holding cost is a cost per unit, rather than a
percentage recalculated from the price.
\path{prob_descuento.json} adds a \(10\%\) discount starting at \(300\)
units. With \texttt{solve-inventory-json}, the selected quantity is
\(Q=300\), with price \(4.5\), setup cost \(120\), holding cost \(300\),
purchasing cost \(5400\), and total cost \(5820\). Relevant cost increases
to \(420\), but the purchasing saving makes this tier preferable. The
candidate table and \texttt{optimum} explain the decision; enumeration of the
tiers is exact within the continuous model.

\section{The newsvendor model}
A single replenishment is chosen for a horizon with no later opportunity to
replenish. For selling price \(v\), purchasing cost \(c\), salvage value
\(s\), and additional shortage penalty \(b\), the marginal costs are
\(C_u=v-c+b\) and \(C_o=c-s\). The critical fractile is
\[
 F(q^*)=\frac{C_u}{C_u+C_o}.
\]
QSB supports normally distributed demand, calculating its quantile and expected
sales, leftovers, shortages, and profit. The probability of covering all demand
is the cycle service level; it is not the expected fraction of demand fulfilled.
In \path{prob_newsboy.json}, mean \(100\), standard deviation \(20\),
\(v=5,c=2,s=1,b=1\), and zero fixed cost give critical fractile \(0.8\),
order quantity \(116.832425\), expected sales \(97.767247\), leftovers
\(19.065178\), shortages \(2.232753\), and profit \(272.003808\).
The normal quantile uses approximate arithmetic; the \texttt{closedForm}
label describes the method, rather than an exact decimal quantity or a guarantee
for nonnormal demand.

\section{Lot sizing}
Ending inventory satisfies \(I_t=I_{t-1}+q_t-d_t\). QSB minimizes
\[
 \sum_t\left(S_t\mathbf{1}_{q_t>0}+c_tq_t+
 h_t\max(I_t,0)+b_t\max(-I_t,0)\right).
\]
Dynamic programming uses integer demand and quantities, zero initial inventory,
and zero final inventory, allowing backorders at a cost. It does not represent
production capacity or configurable initial stock. Exactness refers to the
bounded state space of this educational implementation: inventories between
minus and plus total demand, and production in each period between zero and
that demand. Computational cost grows with total demand and the horizon.
\path{prob_lotes.json} has three demands of ten units, setup cost \(20\),
unit production cost \(2\), holding cost \(1\), and backorder cost \(10\).
QSB returns production \((20,0,10)\), ending inventory \((10,0,0)\), and
cost \(110\). The policy \((10,20,0)\) also costs \(110\); an optimal
result does not imply uniqueness. The decision table allows the balance and
cost allocation to each period to be checked.

\section{Review with uncertain demand}
The four policies use normal demand per period and constant lead time. With
independent demand increments, protection period \(T\) has mean \(DT\)
and standard deviation \(\sigma\sqrt T\). Safety stock is expressed using
the normal quantile, and expected shortage cost uses the normal loss function.
These are continuous-quantity approximations, marked \texttt{approximate};
they do not solve a general optimal policy for arbitrary distributions or
random lead times.

The four examples share annual demand \(1200\), standard deviation \(100\),
lead time \(0.02\) years, setup cost \(30\), holding cost \(2\), backorder
cost \(20\), and purchasing cost \(5\). Every shortage is backordered;
review cost is \(5\). The continuous order-up-to review file uses a supplied
average order size of \(200\).
\begin{center}
\begin{tabular}{lll}
\toprule
File under \texttt{manual/examples/} & Policy & Main result \\
\midrule
\path{prob_qr.json} & Fixed quantity \((Q,r)\) & \(Q=194.866241,r=54.242691\) \\
\path{prob_sS.json} & Continuous order-up-to & \(r=54.095104,S=254.095104\) \\
\path{prob_periodica.json} & Fixed interval & \(T=0.170783,S=321.459989\) \\
\path{prob_opcional.json} & Optional periodic & \(T=0.064550,s=163.951339\) \\
\bottomrule
\end{tabular}
\end{center}
Optional review also returns \(S=353.687998\). Under fixed-quantity continuous
review, safety stock is \(30.242691\), service level \(0.983761\), and
relevant cost \(450.217864\). Under fixed-interval periodic review,
protection includes the review interval plus lead time, \(0.190783\) years:
confusing these intervals artificially reduces the stock required. Backorder
and lost-sales fractions, their costs, and distributions are validated in the
core.

\section{macOS workflow and traceability}
Open the normalized JSON or select an inventory variant in the model creator.
Enter demand, costs, and the variant-specific parameters; validate before
running. Keep units visible and review warnings. In the results, EOQ shows
costs and a comparison with the supplied lot; discounts show candidates; the
newsvendor model shows service and demand outcomes; lot sizing shows quantities
and costs by period; stochastic review shows levels, safety stock, and costs.
The analysis/JSON selector retains the core document for export without
transcribing numbers from charts. Quantities and formulas come from QSBCore
in both the CLI and macOS.

All files in this chapter are validated with \texttt{validate-inventory-json}
and solved with \texttt{solve-inventory-json}; adding
\texttt{--backend validate} validates without solving. The verifier
\path{manual/verify_examples.py} checks the expected results in
\path{manual/examples/manifest.json}; it requires no proprietary reference
material.

\section{Exercises}
\begin{enumerate}
\item Compare integer orders of 189 and 190 in EOQ and explain why comparing
costs is preferable to rounding blindly.
\item Double the lead time in the stochastic example: compare changes in
protected demand and safety stock. Do both double?
\item Lower the backorder cost in the lot-sizing example and check whether
negative inventories appear. Interpret backorders operationally.
\end{enumerate}

\chapter{Scheduling and sequencing}
\label{chap:calendarizacion}

\section{Problem and assumptions}
Scheduling assigns operations to machines and times. QSB minimizes the
\emph{makespan}, the time at which the last operation finishes
\cite{pinedo2016scheduling}. Operations cannot be interrupted, a machine
processes one operation at a time, and each job follows its route. Durations
are integers in the specified unit. Setup times, working calendars, and
random failures are not included.
Due-date and weight fields are preserved, but the native objective remains
makespan; it does not optimize weighted tardiness.

\section{Flow shop: a common route}
In a flow shop, all jobs visit the machines in the same order.
The native solver searches for a permutation shared by all machines. For
a given sequence, an operation starts when both the job and machine are
available; its finish is its start plus its duration. Initial machine
availability times are respected when provided.

In \path{manual/examples/plan_flowShop.json}, there are two jobs and two
machines, initially available at zero:
\begin{center}
\begin{tabular}{lrr}\toprule
Job & M1 & M2 \\
\midrule J1 & 2 & 3 \\ J2 & 3 & 1 \\
\bottomrule\end{tabular}
\end{center}
\begin{lstlisting}[language=bash]
swift run qsb validate-scheduling-json manual/examples/plan_flowShop.json
swift run qsb solve-scheduling-json manual/examples/plan_flowShop.json --backend native
\end{lstlisting}
The sequence J1, J2 has makespan 6. On M1, J1 occupies $[0,2]$ and J2 $[2,5]$;
on M2, J1 occupies $[2,5]$ and J2 $[5,6]$. The alternative J2, J1 finishes at 8.
M2's initial idle time is not a missing operation: it is waiting for J1 to
finish. The search is exact within the permutation model, with a limit of
eight jobs and factorial cost; it does not prove optimality over all possible
machine-specific sequences in a general flow shop.

\section{Job shop: individual routes}
In a job shop, each job may follow a different route. The example
\path{plan_jobShop.json} requires J1: M1(2), M2(3), and J2: M2(2), M1(1).
Precedence constraints follow from the order of each job's operations.
\begin{lstlisting}[language=bash]
swift run qsb validate-scheduling-json manual/examples/plan_jobShop.json
swift run qsb solve-scheduling-json manual/examples/plan_jobShop.json --backend native
\end{lstlisting}
QSB obtains makespan 5: J1 uses M1 from 0 to 2 and M2 from 2 to 5; J2 uses
M2 from 0 to 2 and M1 from 2 to 3. J1's total duration is 5 and provides a
lower bound, so the schedule achieves the optimum. In this case,
\texttt{jobSequence} reflects the operation dispatch order and may repeat
names; it must not be interpreted as a permutation of jobs.
Branch and bound with dominance pruning supports up to six jobs and
36 operations. These are educational limits: search time can grow
substantially well before they are reached. Native is deterministic and
marks both variants \texttt{fixtureScale}; Validate does not construct
schedules, and External is unavailable.

\section{Reading the Gantt chart and using macOS}
Import the JSON envelope, whose \texttt{kind} is \texttt{flowShop} or
\texttt{jobShop}. To enter or change operations, use the structured
Definition form or the model JSON. These variants do not use the CPM/PERT
activity editor. They are different problems, and the data for one do not
replace those for the other.
Validate and run with QSB Native. The solution opens Timeline: each row is
a machine, bars are operations, and gaps show idle time.
The horizontal scale lets you examine long schedules. Operation details
identify the job, machine, start, finish, and duration; JSON also contains
\texttt{idleBefore} and the schedules by machine.

Before accepting the result, check that no machine has overlapping bars,
that each job retains its operation order, and that the maximum finish
equals the makespan. A short schedule does not guarantee that due dates
are met; compare them separately if they matter to the actual decision.
The algorithms run in QSBCore, and the Gantt chart presents the same times
exported by the CLI.

\section{Exercises}
\begin{enumerate}
\item Change J2's second operation in the flow shop from 1 to 5.
      Compare both permutations and verify the result.
\item In the job shop, add an initial availability time of 2 for M2.
      Draw the waiting times and check that the solver does not schedule
      operations before that time.
\item Set J1's due date to 4. Explain why changing this field
      does not turn the problem into tardiness minimization.
\end{enumerate}

\chapter{Projects and production planning}
\label{chap:planificacion}

\section{CPM: precedence and the critical path}
CPM operates on an acyclic activity network: an activity starts only after
all its predecessors have finished. It does not allocate limited resources.
The forward pass calculates early starts and finishes; the backward pass
calculates late starts and finishes. Total float is late start minus early
start. A critical activity has zero float.

In \path{manual/examples/plan_cpm.json}, A takes 3 days, B takes 2, and both
precede C, which takes 4. Each activity has a normal cost of 10, a crash
duration one unit shorter, and a crash cost of 15.
\begin{lstlisting}[language=bash]
swift run qsb validate-project-json manual/examples/plan_cpm.json
swift run qsb solve-project-json manual/examples/plan_cpm.json --backend native
swift run qsb solve-cpm-crash-json manual/examples/plan_cpm.json 6 --backend native
\end{lstlisting}
\begin{figure}[htbp]
\centering
\begin{tikzpicture}
\node[qsb activity,line width=1.4pt] (A) at (0,1) {A\\3};
\node[qsb activity] (B) at (0,-1) {B\\2};
\node[qsb activity,line width=1.4pt] (C) at (4,0) {C\\4};
\draw[qsb selected] (A) -- (C);
\draw[qsb alternative] (B) -- (C);
\end{tikzpicture}

\caption{CPM activity network: the number below each name is its duration in days. Thick strokes identify critical activities and precedence; the dashed stroke shows the other precedence.}
\label{fig:cpm}
\end{figure}

The normal project finishes in 7 days. A and C are critical; B has float 1.
The total normal cost is 30. Crashing to a deadline of 6 reduces C from 4
to 3, with incremental cost 5 and total cost 35. Reducing A could also
achieve the same cost; do not assume that the plan is unique.

Crash optimization assumes linear costs between normal and minimum
durations, with slope $(C_{crash}-C_{normal})/(t_{normal}-t_{crash})$ when
reduction is possible. It uses continuous times and an educational LP that
considers paths; the number of paths can grow rapidly. It is not a policy
for hiring resources and does not include indirect project costs.
Deadlines shorter than what is achievable must be treated as infeasible,
never as instructions to violate precedence constraints.

\section{PERT: expected duration and uncertainty}
PERT uses optimistic $a$, most likely $m$, and pessimistic $b$ estimates,
ordered $a\leq m\leq b$. The classical approximation calculates
\[
 E[T]=(a+4m+b)/6,\qquad \operatorname{Var}(T)=((b-a)/6)^2.
\]
In \path{plan_pert.json}, the triples are A(2,3,4), B(1,2,3), and C(3,4,5),
with the same precedence relationships as in the CPM example.
\begin{lstlisting}[language=bash]
swift run qsb solve-project-json manual/examples/plan_pert.json --backend native
\end{lstlisting}
\begin{figure}[htbp]
\centering
\begin{tikzpicture}
\node[qsb activity,line width=1.4pt] (A) at (0,1) {A\\(2,3,4)\\$E[T]=3$};
\node[qsb activity] (B) at (0,-1) {B\\(1,2,3)\\$E[T]=2$};
\node[qsb activity,line width=1.4pt] (C) at (4,0) {C\\(3,4,5)\\$E[T]=4$};
\draw[qsb selected] (A) -- (C);
\draw[qsb alternative] (B) -- (C);
\end{tikzpicture}

\caption{PERT network: each activity shows its optimistic, most likely, and pessimistic triple and its expected duration in days. Thick strokes identify the critical path calculated using those means.}
\label{fig:pert}
\end{figure}

Scheduling with expected durations again gives a duration of 7, with A and
C critical. The reported variance is $2/9$ and the standard deviation is
$0.471405$ days. In this example, there is a single critical path, and
summing its variances assumes independence. The implementation sums the
variances of all activities marked critical; if there are multiple critical
paths, that total is neither the variance of a single path nor that of the
maximum over them. The critical path may change between realizations;
replacing random durations with their means does not calculate the expected
maximum exactly. The \texttt{exact} label describes the calculation using
those expected durations, not an exact probabilistic prediction of the
completion date. QSB does not return a probability of meeting a deadline
or simulate that uncertainty in this module.

\subsection{CPM and PERT in macOS}
Import the project envelope or create a CPM/PERT definition. The activity
editor collects names, predecessors, and durations; CPM adds costs, and
PERT adds the three estimates. Missing predecessors and cycles appear in
Validation. Switching methods preserves data for both methods after
initialization: check the estimates when returning to PERT.
Copy Table/Paste Table lets you transfer activities from a spreadsheet.
Run produces duration, critical activities, early/late bars, float values,
and a timing table. Review JSON for the variance and exact values.
The CLI crash operation is a separate workflow; the normal view does not
imply that a deadline has been optimized. Native and Validate are
available; External does not currently solve project models.

\section{Aggregate planning}
Aggregate planning balances demand across periods through regular
production, overtime, subcontracting, inventory, and backorders.
When capacity depends on workforce, it adds hiring and layoffs.
A basic balance is
\[
 I_t-B_t=I_{t-1}-B_{t-1}+R_t+O_t+S_t-D_t.
\]
The objective sums production, idle capacity, workforce, inventory, and
backorder costs according to the data. QSB normalizes the historical simple,
transportation, and LP methods into a single continuous LP formulation;
the \texttt{method} field preserves the input identity rather than selecting
three different native algorithms.

The example \path{plan_aggregate.json} has two months, demand 4 and 6,
regular capacity 5 each month, regular cost 2 per unit, overtime capacity
up to 2 at cost 5, inventory cost 1, and no backorders or subcontracting.
It does not depend on workforce and starts with no inventory.
\begin{lstlisting}[language=bash]
swift run qsb validate-aggregate-json manual/examples/plan_aggregate.json
swift run qsb solve-aggregate-json manual/examples/plan_aggregate.json --backend native
\end{lstlisting}
The plan produces 5 in each month; it holds 1 at the end of the first and
zero at the end of the second. It costs $10\cdot2+1=21$. Producing exactly
the first month's demand and using one overtime unit in the second would
cost $9\cdot2+5=23$.
Units must be consistent: capacity per worker is not total capacity, and
\texttt{capacityRequirementPerUnit} converts production into capacity.
Vectors must have the same number of periods; \texttt{null} in optional
maximum fields means no upper bound. Workforce and quantities are
continuous, even when the business needs whole people. Do not round a
solution without rechecking balances and capacity.

In macOS, the period editor separates demand, costs, inventory/sources,
and workforce. Copy Table transfers all 20 period columns; Paste Table
validates the range before replacing it. Run shows stacked production and
demand, inventory, backorders, and a plan table. A backorder appears below
zero in the chart, but its model quantity is nonnegative. Native solves
the educational LP; External uses HiGHS for the same formulation if
installed. Validate does not produce a plan. The optimal objective refers
to the costs and bounds entered, not to omitted criteria such as workforce
stability and its social implications.

\section{MRP: requirements explosion}
MRP starts from a master production schedule (MPS), a bill of materials
(BOM), inventory, scheduled receipts, and lead times. It processes parents
before components: a parent's order release generates gross requirements
for its components, multiplied by the quantity used per unit
\cite{silver1998inventory}.
For an item with safety stock $SS$, gross requirement $G_t$, and scheduled
receipt $SR_t$, net requirement is
$\max(0,G_t+SS-I_{t-1}-SR_t)$. A planned receipt covers that requirement,
and its release is shifted backward by the lead time.

\path{plan_mrp.json} has weeks S0--S3. Five units of P must be received in
S3; P uses two units of C per unit. Both have a one-week lead time, zero
initial inventory and safety stock, the LFL rule, and unlimited capacity.
\begin{lstlisting}[language=bash]
swift run qsb validate-mrp-json manual/examples/plan_mrp.json
swift run qsb solve-mrp-json manual/examples/plan_mrp.json --backend native
\end{lstlisting}
P receives 5 in S3 and releases 5 in S2. C needs 10 in S2, receives 10
then, and releases 10 in S1. Projected inventories remain zero. C's
requirements come from P's release, not from the final demand date.
The BOM must be acyclic, and all its identifiers must exist.

The available rules are LFL, EOQ, LUC, LTC, and PPB. LFL receives the net
requirement; EOQ starts from $\sqrt{2DS/H}$ and receives at least the
current requirement. In EOQ, \texttt{lotSizingLookahead} is interpreted as
a lot multiple \(m=\max(1,\text{configured value})\), not as a number of
periods: the quantity is
\(\max(\text{requirement},m\lceil EOQ/m\rceil)\).
If annual holding cost is zero, it receives the net requirement.
LUC compares unit cost; LTC and PPB currently use the same selector, which
minimizes the absolute difference between holding cost and ordering cost
within the configured horizon. These are forward-looking rules, not a
global optimization of multilevel costs. QSB labels the MRP document
\texttt{heuristic}, even when the particular case uses LFL.
Finite capacity generates \texttt{capacityExcess}; it does not level or
reschedule orders. If a release falls before the horizon, QSB accumulates
it in the first bucket using $\max(0,t-L)$; that record requires operational
review and does not prove that the lead time can be met.

In macOS, Item Master collects parameters, Schedules collects quantities
by period, and Bill of Materials collects parent/component relationships.
Select an item to edit or read its schedule; long horizons are grouped for
navigation. Copy/Paste Table always exchanges the complete range for its
section. The solution shows requirements, inventory, receipts, releases,
and capacity excesses by item, with JSON as an audit record. Native and
Validate are available; External is not implemented for MRP.

\section{Exercises}
\begin{enumerate}
\item In CPM, increase B from 2 to 4. Recalculate the critical path and float.
\item In PERT, increase only A's pessimistic estimate. Compare duration and
      variance; explain which assumptions limit a probabilistic interpretation.
\item Raise aggregate inventory cost from 1 to 4. Compare early production
      with overtime.
\item Add initial inventory of 4 units of C in MRP and check its net
      requirement. Then limit C's capacity in S2 to 5 and explain the excess.
\end{enumerate}

\chapter{Facilities and line balancing}
\label{chap:instalaciones}

\section{Locating a facility}
Continuous location minimizes weighted distance to existing points:

\[\min_{x,y}\sum_i w_i d((x,y),(x_i,y_i)).\] Weights represent interactions
multiplied by unit cost, not probabilities. QSB supports one new facility
and minimization; it does not solve discrete warehouse selection, multiple
new facilities, or problems with geographic barriers.

The distance measure determines the algorithm. For rectilinear distance,
$d=|x-x_i|+|y-y_i|$, weighted medians are used coordinate by coordinate.
For squared Euclidean distance, $d=(x-x_i)^2+(y-y_i)^2$, the weighted
centroid is used. For ordinary Euclidean distance, Weiszfeld iteration is
used and the result is approximate. Do not change the measure merely to
obtain visually more attractive coordinates: it changes the objective and
its units.

\path{manual/examples/plan_location.json} places F1 at $(0,0)$ with weight
1 and F2 at $(4,0)$ with weight 3, and uses squared Euclidean distance.
The directed matrix contains these values from the existing facilities
to the new one; there are no other interactions. The effective weight
sums both interaction directions between each existing point and the new
one: recording the same flow in both directions doubles the weight.
\begin{lstlisting}[language=bash]
swift run qsb validate-facilities-json manual/examples/plan_location.json
swift run qsb solve-facilities-json manual/examples/plan_location.json --backend native
\end{lstlisting}
The center is $(3,0)$ and the objective is $1\cdot9+3\cdot1=12$.
In macOS, Definition lets you edit roles, IDs, coordinates, distance, and
the matrix. New coordinates may be left blank; existing coordinates are
required. IDs are interpreted within the existing/new role, so both can
use ID 1. The solution displays points and weighted connections and
details each contribution. Use JSON to preserve the complete matrix.

\section{Balancing a line}
With indivisible tasks of duration $t_i$, precedence constraints, and fixed
cycle time $C$, the objective minimizes the number of stations. A station
cannot exceed the cycle time, and a predecessor's station cannot follow
its successor's station; predecessor and successor tasks may share a
station if performed in order. The lower bound
$\lceil\sum_i t_i/C\rceil$ ignores incompatibilities and precedence, so it
is not always attainable.

\path{plan_line.json} contains A(2), B(3), C(4), with A and B preceding C,
and a cycle time of 5 minutes.
\begin{lstlisting}[language=bash]
swift run qsb validate-facilities-json manual/examples/plan_line.json
swift run qsb solve-facilities-json manual/examples/plan_line.json --backend native
\end{lstlisting}
A and B are assigned to station 1 and C to station 2. There are 9 minutes
of work and 10 minutes of capacity per cycle: efficiency is
$9/(2\cdot5)=0.9$ and balance delay is $1-0.9=0.1$. Idle time is 0 and
1 minute, respectively. The lower bound is 2 and is attained, proving
optimality for the model.
Exact bitmask dynamic programming supports at most 24 tasks; a task
longer than the cycle time makes the problem infeasible.
The \texttt{isolated} field is preserved for compatibility, but the current
native solver does not require a separate station for marked tasks. Do
not use it as an effective safety or incompatibility constraint.

The macOS editor supports tasks, durations, immediate successors, and
flags; Copy/Paste Table transfers the table. Deleting a referenced task
leaves the reference visible so that it can be repaired. The solution
shows station workload and idle time, assignments, precedence context,
efficiency, and balance delay. Efficiency does not measure total batch
production time or include breakdowns or product changeovers.

\section{Facility layout}
Layout evaluates departments positioned on a grid. Each department
occupies the union of one or more cell rectangles; a repeated cell within
the same department is counted once. Rectangle coordinates are inclusive
and start at 1. The load--distance objective uses centroids:
\[
 Z=\sum_{i,j}w_{ij}(|r_i-r_j|+|c_i-c_j|).
\]
The matrix is directed: entering the same weight in both directions
doubles the contribution. The centroid approximates physical travel;
it does not model doors, mandatory aisles, or obstacles.

The example \path{plan_layout.json} has one row and three cells A, B, C,
all with area one and movable. Its only weights are A--B=1, A--C=10,
and B--C=1.
\begin{lstlisting}[language=bash]
swift run qsb validate-facilities-json manual/examples/plan_layout.json
swift run qsb solve-facilities-json manual/examples/plan_layout.json --backend native
swift run qsb solve-facilities-json manual/examples/plan_layout.json \
  --backend native --layout-strategy pairwise-swap
\end{lstlisting}
Initial evaluation gives $1+20+1=22$ and does not change the arrangement.
Pairwise swaps A and B: B, A, C costs $1+10+2=13$, an improvement of 9.
The search swaps only movable departments of equal area when doing so
improves the objective. It terminates at a local optimum in that
neighborhood; it does not guarantee a global optimum or automatically
construct arbitrary geometry. Fixed areas are not swapped. Overlaps
between movable areas are errors; overlaps involving fixed areas may
produce warnings and must be reviewed.

In Definition, you edit the grid, departments, fixed flags, rectangles,
and flow/cost matrix. Select a department to access its rectangles;
the transferred table includes IDs for all of them. In large matrices,
sources and destination groups are selected, although Copy/Paste Table
transfers the complete matrix. Choose Initial or Pairwise before Run.
The Layout view lets you adjust scale, grid, and flows, and examine
centroids, swaps, and contributions. For large grids, the overview draws
fewer lines; exact ranges remain available in details and JSON.

\section{Backend, validation, and decision}
All three families support Native and Validate, with no operational
External backend. Validate reports invalid data without calculating a
location, layout, or line balance. Documents preserve algorithm and
exactness: centroid/median location is closed-form, Weiszfeld is
approximate, line balancing is exact at educational scale, and Pairwise
is heuristic. Evaluating an initial layout determines only its cost,
not its optimality.
The examples use the facilities envelope with \texttt{kind} equal to
\texttt{location}, \texttt{lineBalancing}, or \texttt{layout}. Do not confuse
that envelope with the individual JSON formats accepted by commands
specific to each subfamily. Validate the same file you will solve and
export the model and solution to retain a traceable record of accepted data.

\section{Exercises}
\begin{enumerate}
\item Change location to rectilinear distance. Explain the new coordinate
      using weighted medians and the new units of the objective.
\item Reduce the line's cycle time to 4. Determine the minimum number of
      stations and check the impact on efficiency.
\item Fix A in the layout. Solve with Initial and Pairwise and compare
      permitted moves. Explain why a search with no improvement does not
      prove optimality over all possible layouts.
\end{enumerate}


\part{Uncertainty and quality}
\chapter{Queuing models}
\label{ch:colas}

\section{From flow to waiting}
Queuing theory relates variability, capacity, and congestion
\cite{gross2008fundamentals}. In Kendall notation, M denotes Poisson arrivals
or exponential times, depending on its position. Service rate \(\mu\) is
per server; mean service time is \(1/\mu\). All rates share a time unit.
QSB implements M/M/1 and a finite queue with multiple servers using a
birth-and-death model.

\section{M/M/1 in steady state}
The assumptions are independent arrivals, exponential service, one server,
and unlimited population and waiting space. Stability requires
\(\lambda<\mu\). With \(\rho=\lambda/\mu\):
\[
 P_0=1-\rho,\quad L=\frac{\rho}{1-\rho},\quad
 L_q=\frac{\rho^2}{1-\rho},\quad
 W=\frac1{\mu-\lambda},\quad W_q=\frac{\lambda}{\mu(\mu-\lambda)}.
\]
\(L\) includes customers being served; \(L_q\) includes only those waiting.
Little's law gives \(L=\lambda W\) and \(L_q=\lambda W_q\).
The backend reports \texttt{mm1ClosedForm}/\texttt{closedForm}. It does not
calculate the distribution or percentiles of an individual waiting time.

\subsection{Example: a service counter}
\path{manual/examples/prob_cola.json} specifies four arrivals/hour and six
services/hour. Run from the root:
\begin{lstlisting}
.build/debug/qsb validate-queuing-json manual/examples/prob_cola.json
.build/debug/qsb solve-queuing-json manual/examples/prob_cola.json
\end{lstlisting}
Utilization is \(2/3\), \(P_0=1/3\), \(L=2\), \(L_q=4/3\),
\(W=0.5\) hours, and \(W_q=1/3\) hours: thirty minutes in the system,
twenty waiting, and ten receiving service. Being busy two thirds of the time
does not mean that two thirds of customers wait exactly twenty minutes.
If arrivals reach six/hour, the model no longer has these steady-state means,
and QSB returns a stability diagnostic.

\section{Multiple servers and finite capacity}
With \(c\) servers and \(b\) waiting places, system capacity is \(K=c+b\).
State \(n\) has arrival rate \(\lambda\) when \(n<K\), and zero when
full; its departure rate is \(\min(n,c)\mu\).
The probabilities are normalized after being obtained from
\[
 p_n=p_0\prod_{i=1}^n\frac{\lambda}{\min(i,c)\mu}.
\]
Blocking probability is \(p_K\), the effective arrival rate
is \(\lambda_e=\lambda(1-p_K)\), and Little's law uses \(\lambda_e\),
rather than the offered rate. Because it is finite, this model can have a
stationary distribution even when offered demand exceeds its service capacity;
this produces blocking, rather than extra capacity.

QSB requires exponential arrivals in batches of one and an infinite customer population.
It uses mean service time to obtain the rate per server; nonexponential labels
produce warnings because their variability is not modeled. The document always
retains \texttt{approximate} for this route. A case with normal service times
must not be presented as an exact M/G/c/K solution.

\subsection{Example: two counters and three waiting places}
\path{prob_cola_finita.json} retains four arrivals/hour and a mean service
time of ten minutes, with two servers and three waiting places. Validate and
solve using the same commands with the file changed. The results are:
\begin{center}
\begin{tabular}{lr}
\toprule
Metric & Value \\
\midrule
Total capacity & 5 \\
Blocking probability & 0.004124 \\
Effective arrivals per hour & 3.983505 \\
Average utilization per server & 0.331959 \\
Mean number in the system & 0.738144 \\
Mean number waiting & 0.074227 \\
Time in the system (hours) & 0.185300 \\
Waiting time (hours) & 0.018634 \\
\bottomrule
\end{tabular}
\end{center}
Mean waiting time is approximately \(1.118012\) minutes among admitted
customers. The total utilization of both servers corresponds to \(0.663918\)
busy servers, rather than the individual utilization percentage.

\section{Costs, interface, and limitations}
Optional costs distinguish busy and idle servers, waiting customers, and
customers in service; the finite queue adds rejection and waiting places.
State costs are weighted by expected occupancy; rejection cost is multiplied
by the rate of blocked customers. The total is a cost per time unit, rather
than per customer.

In macOS, the creator offers both variants. Enter rates for M/M/1, or mean
times, servers, and waiting space for a finite queue; validate and solve. The
metrics panel shows utilization, effective flow, blocking, numbers, and times.
The finite variant also shows state probabilities, highlighting the full state;
optional costs are itemized. Check the time unit and assumptions alongside the
results. Normalized JSON is the alternative for inspection and export.

Use \path{solve-queuing-json} to accept normalized input.

The historical commands \path{solve-mm1-json} and
\path{solve-finite-queue-json} accept a legacy file despite their suffix:
the suffix indicates JSON output. Always use the generic route to exchange
normalized inputs. Native and validation backends are available; no external
solver is installed for this family.

Priorities, customer abandonment due to impatience, finite customer populations,
general batch arrivals, and rate changes by time of day are not modeled.
The simulation in the Markov and simulation chapter allows a different class
of assumptions to be explored, but its estimates do not replace these
steady-state formulas.

\section{Exercises}
\begin{enumerate}
\item Change \(\lambda\) to 5.5 in M/M/1 and compare utilization and waiting.
Explain why the growth in waiting time is nonlinear.
\item Remove the waiting places from the finite queue. Check that blocking
occurs without a positive queue length, and use \(\lambda_e\) to verify
Little's law.
\item Compare adding a server with adding five waiting places: which changes
service capacity, and which changes the space for customers?
\end{enumerate}

\chapter{Forecasting and regression}
\label{ch:pronosticos}

\section{A forecast is an operational hypothesis}
Forecasts inform inventory, capacity, and planning. A straight line that fits
the past well does not guarantee future accuracy
\cite{hyndman2021forecasting}. QSB implements linear trend, moving average,
simple exponential smoothing, multiplicative seasonal decomposition, and
least-squares linear regression. The JSON is a request: it contains the model,
method, horizon, and parameters. This preserves the chosen method alongside
the data. Native and validation backends are available; external integration
is a future extension.

\section{Data and errors}
Observations are processed in row order; labels do not reorder dates or fill
missing periods. A monthly series must have consistent monthly spacing. QSB
returns fitted values, residuals \(e_t=y_t-\hat y_t\), MAD, MSE, and MAPE:
\[
 \mathrm{MAD}=\frac1m\sum|e_t|,\quad
 \mathrm{MSE}=\frac1m\sum e_t^2,\quad
 \mathrm{MAPE}=\frac{100}{m_+}\sum_{|y_t|>10^{-12}}|e_t/y_t|.
\]
\(m\) counts the available fitted rows, rather than necessarily the entire
history; \(m_+\) excludes actual values whose magnitude is no greater than
\(10^{-12}\). If all are excluded, MAPE is unavailable. MAD retains the
series' units; MSE has squared units; MAPE is reported as a percentage.
Comparing methods requires considering that moving average and smoothing omit
different initial rows.

\section{Linear trend: a complete example}
The model \(y_t=a+bt+e_t\) estimates \(a,b\) by least squares with indices
\(t=1,\ldots,n\). \path{manual/examples/prob_pronostico.json} contains
\((12,14,16,18,20,22)\), month as its unit, and a horizon of two. From the
root:
\begin{lstlisting}
.build/debug/qsb validate-forecast-json manual/examples/prob_pronostico.json
.build/debug/qsb solve-forecast-json manual/examples/prob_pronostico.json
\end{lstlisting}
QSB obtains \(a=10,b=2\), forecasts of \(24\) and \(26\), and zero for
all three error measures. These artificial data were chosen to check the
formula and indices: zero error is not evidence about real demand. The result
retains the request, the \texttt{linearTrend} algorithm, and \texttt{exact}
exactness, which describes the method's calculation rather than exact knowledge
of the future.

\section{Moving average and smoothing}
For window \(k\), the one-step forecast is
\[
 \hat y_{t+1}=\frac1k\sum_{j=0}^{k-1}y_{t-j}.
\]
QSB extends horizons by using earlier forecasts within the window when actual
values run out. The file \path{prob_movil.json} selects a window of three
for the same series. The results are \(20\) and \(20.666667\): the second
average uses \((20,22,20)\). MAD is \(4\), MSE \(16\), and MAPE
\(20.134680\%\), calculated over the final three rows. This lag is expected
when the series has a trend.

Simple smoothing applies
\[
 F_{t+1}=\alpha y_t+(1-\alpha)F_t,\qquad 0<\alpha\le1.
\]
The initial forecast is set to the first observation. In
\path{prob_suavizado.json}, \(\alpha=0.3\) produces \(18.117660\) in
both future periods. The method has no trend component; it does not
automatically extrapolate growth. Its MAD is \(4.078440\), MSE
\(18.233818\), and MAPE \(22.081810\%\). A large alpha reacts quickly
and may follow more noise; a small one smooths and may lag.
Both examples run with \texttt{solve-forecast-json} and are validated with
\texttt{validate-forecast-json}.

\section{Multiplicative seasonality}
The method combines the estimated trend with seasonal factors:
\(\hat y_t=(a+bt)s_j\). QSB calculates actual/trend ratios, averages them
by seasonal position, and normalizes the factors to have a mean of one.
This is an educational decomposition, rather than Holt--Winters or an ARIMA
model. Seasonal length and observation order must correspond to the phenomenon;
at least two complete cycles and values/trend compatible with multiplicative
division are required.
\path{prob_estacional.json} uses length two for the previous linear series.
Both factors are one; the forecast is again \(24,26\). The case checks that
specifying a season does not invent oscillations when data coincide with the
line. To study actual seasonality, replace the series with several oscillating
cycles and review residuals before interpreting the factors.

\section{Explanatory regression}
The model is \(y_i=\beta_0+\sum_j\beta_jx_{ij}+e_i\), with an intercept.
Predictors must have linearly independent columns; a constant column duplicates
the intercept. The core detects rank deficiency.
\path{prob_regresion.json} contains activity \(1,2,3,4,5\) and consumption
\(5,7,9,11,13\). With \texttt{solve-forecast-json}, QSB returns intercept
\(3\), activity coefficient \(2\), zero SSE, and \(R^2=1\).
The request contract's horizon must also be positive for regression, although
its output fits the supplied observations and does not generate future predictor
rows. For a new activity level, the estimated equation allows a prediction to
be calculated, without QSB reporting an interval for it.

A coefficient is a conditional association, rather than causality. \(R^2\)
measures within-sample fit; it does not identify overfitting or guarantee
extrapolation. This workflow does not report coefficient standard errors,
hypothesis tests, or prediction intervals. Small original datasets allow the
contract to be checked, but do not assess generalization ability.

\section{Working in macOS and evaluation}
In the forecasting editor, select the method, parameters, and horizon; enter
historical rows in temporal order. Values can be pasted into one column, or
labels and tab-separated values into two columns. Paste validation is atomic:
a malformed row does not partially replace the history. Regression uses
variable names and explanatory rows. Validate, solve, and review actual,
fitted, and future curves; these curves and the residual chart share observation
order. Visibility and horizontal-scale controls help with long series. The
summary shows parameters and metrics; the JSON selector preserves the request
and result. Applying JSON updates the draft only when decoding succeeds.

To evaluate an operational forecast, reserve the final periods, fit using the
earlier history, and compare against the held-out data. The errors QSB reports
are the fits from the executed calculation; they do not automatically constitute
an out-of-sample test. Do not select parameters by repeatedly examining the
same held-out set. Document units, window, alpha, season, and horizon when
delivering the analysis.

\section{Exercises}
\begin{enumerate}
\item Compare windows of two and four for the growing series. Explain the lag
and check which rows contribute to MAD.
\item Replace one observation with zero and check the definition of MAPE.
Compare its suitability with MAD.
\item Enter a second predictor equal to twice the first and observe the rank
diagnostic. Explain why more columns do not always add information.
\end{enumerate}

\chapter{Markov processes and simulation}
\label{ch:markov-simulacion}

\section{Two ways to describe uncertainty}
A Markov chain propagates probabilities among states; a simulation constructs
realizations of events and estimates measures through repetition. The former
can yield a stationary distribution through equations; the latter supports
more detailed networks and rules, at the cost of sampling error. QSB maintains
separate normalized models and results for the two procedures
\cite{ross2019introduction,law2015simulation}.

\section{Finite chains and the row convention}
\(P_{ij}\) is the probability of moving from state \(i\) to \(j\) in one
period. Each row sums to one. The distribution is a row vector:
\[
 s_{t+1}=s_tP,\qquad \pi P=\pi,\qquad \sum_i\pi_i=1.
\]
Transitions are assumed time homogeneous, and the next state depends on the
current one without additional information from the past. Costs \(c_i\) are
per state per period; their expectation is \(s_tc\). They are not transition
costs.

\subsection{Example: machine availability}
\path{manual/examples/prob_markov.json} defines Operativa (operational) and
Parada (stopped), costs \((0,100)\), initial distribution \((1,0)\), a
horizon of three, and
\[
 P=\begin{pmatrix}0.9&0.1\\0.4&0.6\end{pmatrix}.
\]
From the repository root:
\begin{lstlisting}
.build/debug/qsb validate-markov-json manual/examples/prob_markov.json
.build/debug/qsb solve-markov-json manual/examples/prob_markov.json
\end{lstlisting}
The results are:
\begin{center}
\begin{tabular}{rrrr}
\toprule
Period & Operativa & Parada & Expected cost \\
\midrule
0 & 1 & 0 & 0 \\
1 & 0.9 & 0.1 & 10 \\
2 & 0.85 & 0.15 & 15 \\
3 & 0.825 & 0.175 & 17.5 \\
\bottomrule
\end{tabular}
\end{center}
The stationary distribution is \((0.8,0.2)\), with expected cost \(20\)
per period. This is a long-run mean, rather than a cumulative cost over three
periods. This chain is irreducible and aperiodic, so its transient distributions
converge to it. In general, stationary uniqueness and convergence from every
initial distribution are not equivalent: a periodic chain may oscillate.

QSB solves a stationary linear system and propagates matrices; it reports
\texttt{exact} as the algorithmic character, using floating-point arithmetic.
The method requires a unique stationary distribution and returns an error if
its system is singular. It does not choose an arbitrary mixture for chains
with several closed classes. When no initial distribution is supplied, it
does not invent one: transient output is empty. Validating rows and dimensions
does not guarantee that the stationary system is solvable.

In macOS, the editor supports names, costs, optional initial probabilities,
the transition matrix, and the horizon. Check row orientation before solving;
the result shows the stationary distribution, propagation, costs, and
assumptions. Editing states keeps dimensions consistent; invalid text remains
visible until corrected. JSON remains available for interchange.

\section{Discrete-event simulation}
A network contains entity sources, queues, servers, and gates. Routes include
probabilities and transfer times; service rules may depend on entity type.
QSB supports constant, exponential, normal, uniform, and triangular
distributions. Negative normal samples are clipped to zero; they are not
resampled. This clipping changes the mean and variability relative to an
unclipped normal distribution. Parameters represent distribution times or
values, rather than rates per time unit. For example, an exponential mean of
\(0.25\) hours corresponds to four arrivals/hour.

The educational backend models finite FIFO queues, alternative servers,
probabilistic routing, transfers, and assembly synchronization. Its outputs
are \texttt{approximate} estimates, even when a fixed seed makes them
reproducible. There is no warm-up deletion, interactive event trace, or
guarantee of steady state. No external simulation backend is available either.

\subsection{Example: a simulated service counter}
\path{prob_simulacion.json} connects source Cliente (customer), queue Espera
(waiting) with twenty places, and server Atención (service). Each arrival
brings one entity, interarrival time is exponential with mean \(0.25\)
hours, and mean exponential service time is \(1/6\) hours. Transfers take
zero time.
\begin{lstlisting}
.build/debug/qsb validate-simulation-json manual/examples/prob_simulacion.json
.build/debug/qsb solve-simulation-json manual/examples/prob_simulacion.json
.build/debug/qsb solve-simulation-json manual/examples/prob_simulacion.json --replications 3
\end{lstlisting}
The default run lasts \(1000\) hours, with seed \(1\), and no warm-up
removed. It generates \(4005\) entities, completes \(4004\), rejects
none, and estimates mean queue length \(1.202238\), maximum \(15\),
and utilization \(0.664470\). The uncompleted customer may be in the queue
or in service at termination: it is not necessarily rejected.

With three independently seeded replications, mean queue length is
\(1.229240\), with approximate normal \(95\%\) interval
\([1.146668,1.311813]\); utilization is \(0.659950\), with interval
\([0.638179,0.681722]\). Reported aggregate counters are rounded means:
\(4006\) generated and \(4005\) completed, rather than the sum of three
runs. Maxima, however, retain the maximum across replications. The base seed,
horizon, and number of replications appear in the output. The CLI exposes
replications but no options for choosing a seed or horizon in these commands;
those values come from the core's default options.

The interval formula uses \(\bar x\pm1.96s/\sqrt R\). With few
replications, the normal approximation may be poor; increasing \(R\) reduces
sampling error, but does not correct startup bias or an incorrect distribution.
Do not interpret the interval as a guarantee that the model describes reality.
The analytical M/M/1 in the queuing chapter gives \(L_q=4/3\) and
\(\rho=2/3\). The simulation has a finite queue and finite horizon:
closeness is useful for checking units, but does not require exact equality.
Three replications do not justify inferring a failure because an interval does
not contain the analytical value, either.

\section{macOS workflow}
In the simulation definition, complete \texttt{Components}, \texttt{Routes}, and \texttt{Service Rules}. Enter the parameters of each
distribution separately. Names must match route references and entity types.
Complete tables can be exchanged with spreadsheets; pasting checks each
section before replacing it. In the Run panel, select between one and one
hundred replications, validate, and run. The interface shows computation status
and prevents duplicate runs; changing the model invalidates a pending result.
This does not stop a calculation already started in the core.

The results panel shows entities generated/completed, completion rate, horizon,
seed, queue, and utilization, with intervals when available. Consult tables
and JSON as well as charts. Keep the network, rules, seed, horizon, and
replications when comparing scenarios: changing any of these may change the
interpretation of the difference.

\section{Exercises}
\begin{enumerate}
\item Remove the initial distribution from the Markov case and check that
only stationary analysis remains. Then change the matrix to the identity and
observe the singular-system error.
\item Double the simulated mean service time and compare completions,
rejections, and utilization. Explain why queue capacity is not equivalent to
service capacity.
\item Repeat the case with ten replications and compare interval widths.
Distinguish variation between runs from differences between model assumptions.
\end{enumerate}

\chapter{Quality control and acceptance sampling}
\label{ch:calidad}

\section{Monitoring a process or accepting a lot}
Statistical control looks for signals of change in a process; acceptance
sampling makes decisions about lots using samples. A stable process can produce
defects, and accepting a lot does not certify that all its units conform.
Both procedures require consistent definitions and sampling
\cite{montgomery2020quality}.

QSB supports c, p, and \(\bar X\)-R charts, Pareto charts, and normal
probability plots; it also evaluates single and double binomial acceptance
plans. The native backend is deterministic, and the validation backend does
not solve. No external backend is available for these families.

\section{c and p charts}
The c chart counts defects in an inspection unit of constant size. An item
may have more than one defect. Under a Poisson count assumption,
\[
 CL=\bar c,\qquad LCL=\max(0,\bar c-3\sqrt{\bar c}),\quad
 UCL=\bar c+3\sqrt{\bar c}.
\]
\path{manual/examples/prob_calidad.json} contains counts
\((2,3,4,3,2,4,3,3,2,4)\). From the root:
\begin{lstlisting}
.build/debug/qsb validate-quality-json manual/examples/prob_calidad.json
.build/debug/qsb solve-quality-json manual/examples/prob_calidad.json
\end{lstlisting}
The center line is \(3\), the lower limit zero, and the upper limit
\(8.196152\). There are no points outside the limits or signals from the
implemented rules. This is an educational check of the formula; ten counts
are insufficient to establish an industrial process as stable. Inspection
units of varying sizes require another approach; QSB does not implement a u
chart in this route.

The p chart uses the proportion of defective units, with size \(n_i\) for
each subgroup. Its weighted center is
\[
 \bar p=\frac{\sum_i n_ip_i}{\sum_i n_i},\qquad
 CL_i=\bar p,\quad CL_i\pm3\sqrt{\bar p(1-\bar p)/n_i},
\]
with limits clipped to \([0,1]\).

In \path{prob_p.json}, sizes
\((100,100,200)\) and proportions \((0.02,0.04,0.03)\) produce
\(\bar p=0.03\); upper limits are \(0.081176\) for the first two
subgroups and \(0.066187\) for the third. The lower limit is zero.
Input uses fractions: \(0.03\) means \(3\%\), rather than \(0.03\%\).
Three-sigma limits are a statistical approximation to the data distribution,
even though QSB labels the application of the formulas as \texttt{exact}.

\section{Mean and range charts}
Rational subgroups combine measurements taken under comparable conditions.
With \(\bar{\bar X}\) as the mean of subgroup means and \(\bar R\)
as the mean of ranges,
\[
 LCL_{\bar X},UCL_{\bar X}=\bar{\bar X}\mp A_2\bar R,\qquad
 LCL_R=D_3\bar R,\quad UCL_R=D_4\bar R.
\]
QSB uses standard constants for uniform sizes between two and ten. Review
variation in R before interpreting shifts in means.
\path{prob_xr.json} contains subgroups \((9,10,11)\), \((10,11,12)\),
and \((8,9,10)\). The results are \(\bar{\bar X}=10\), \(\bar R=2\),
mean limits \([7.954,12.046]\), and range limits \([0,5.148]\), with no
points outside the limits. Control limits describe process variability; they
are not design tolerances. QSB does not calculate Cp/Cpk indices in this
workflow.

\section{Pareto and normal probability}
Pareto aggregates categories, sorts them by descending frequency, and reports
shares and cumulative shares. \path{prob_pareto.json} has rows \((4,2,1)\)
and \((2,1,0)\) for Rayas (scratches), Manchas (stains), and Roturas
(breakages). Their totals are \((6,3,1)\), the overall total is \(10\),
shares are \((0.6,0.3,0.1)\), and cumulative shares are \((0.6,0.9,1)\).
Fields named \texttt{percentage} in JSON store fractions, rather than values
from zero to one hundred. Prioritizing Rayas by frequency does not establish
that it is the category with the greatest cost or risk.

The normal plot sorts measurements and assigns Blom plotting positions
\(p_i=(i-3/8)/(n+1/4)\); it uses \(z_i=\Phi^{-1}(p_i)\), fitting
measurement against normal score. \path{prob_normal.json} contains
\((8,9,10,11,12)\); mean \(10\), sample standard deviation \(1.581139\),
intercept \(10\), slope \(1.742919\), and correlation \(0.997902\).
It is labeled \texttt{approximate} because of the approximate normal quantile.
High correlation with five points does not demonstrate normality; this is not
a test with a p-value. These four examples run using the same quality commands
with the file changed.

\section{Signals and the quality interface}
QSB flags points beyond three sigma and fixed Western Electric rules: two of
three beyond two sigma on the same side, four of five beyond one sigma on the
same side, and eight on the same side of the center. Points already flagged
beyond three sigma are excluded from the counts for the two-of-three and
four-of-five zone rules. There is no rule configuration or normalized metadata
for cause, action, and comment. Signal indices identify observations to
investigate, rather than automatic cause diagnoses. Clipped attribute-chart
limits affect the scale used by this implementation's rules: retain the output
and review each signal alongside the chart.

In macOS, select the variant in the native editor and complete its table:
counts/proportions, subgroups, categories, or measurements. Changing variants
preserves inactive tables during the session, but JSON exports only the active
one. Paste complete tables to replace them atomically; incomplete imported
data remain as invalid cells. After validation and solving, charts show centers,
limits, and signals; Pareto shows frequency and cumulative share; the normal
plot shows points and a fitted line. Consult tables and JSON for the precise
values of the calculation, and preserve acquisition order.

\section{Single acceptance plan}
For sample size \(n\) and maximum acceptable number of defective units
\(c\), the lot is accepted if \(X\le c\). Under independent binomial
sampling with proportion \(p\), the operating characteristic curve is
\[
 P_a(p)=\sum_{x=0}^c{n\choose x}p^x(1-p)^{n-x}.
\]
AQL is the acceptable quality level and RQL the rejectable quality level.
Actual producer's risk is \(\alpha=1-P_a(AQL)\); consumer's risk is
\(\beta=P_a(RQL)\). Entering nominal risks does not alter an already
supplied plan to satisfy them.

\path{prob_muestreo.json} sets \(n=20,c=1,AQL=0.01,RQL=0.15\), nominal
risks \(0.05,0.10\), and lot size \(500\).
\begin{lstlisting}
.build/debug/qsb validate-acceptance-json manual/examples/prob_muestreo.json
.build/debug/qsb solve-acceptance-json manual/examples/prob_muestreo.json
\end{lstlisting}
QSB returns \(P_a(AQL)=0.983141\), producer's risk \(0.016859\), and
consumer's risk \(0.175558\). The plan exceeds the nominal consumer's risk:
it must not be described as satisfying \(\beta=0.10\). Output contains
101 curve points from \(0\) to \(100\%\) defective, as well as measures
at AQL and RQL. This grid does not replace continuous evaluation when optimizing
a risk between points.

Average sample number ASN is \(20\). Under rectifying inspection of rejected
lots and removal of inspected defective units,
\[
 ATI=n+(N-n)(1-P_a),\qquad AOQ=pP_a(N-n)/N.
\]
At AQL, ATI is \(28.092482\) and AOQ \(0.009438\), that is
\(0.943815\%\). If rejected lots are returned without full inspection,
these measures do not describe the actual procedure. QSB evaluates binomial
probabilities exactly within the assumption, but sampling a large fraction
of the lot without replacement requires considering the hypergeometric
distribution; the core provides that probability, rather than a selector in
this command.

\section{Double plans and design}
In a double plan, \(n_1\) units are inspected. If the number of defective
units is \(x_1\le c_1\), the lot is accepted; if \(x_1\ge r_1\), it is
rejected; in between, \(n_2\) \emph{additional} units are inspected. Final
acceptance occurs if \(x_1+x_2\le c_2\). Confusing \(n_2\) with cumulative
sample size changes the plan. Its acceptance probability is
\[
 P_a=P(X_1\le c_1)+
 \sum_{x=c_1+1}^{r_1-1}P(X_1=x)P(X_2\le c_2-x).
\]
The file \path{prob_doble.json} uses \(n_1=10,c_1=0,r_1=3\),
\(n_2=10,c_2=1\), with the same quality levels and lot size as before.
Changing the file in \texttt{solve-acceptance-json} gives producer's risk
\(0.013001\), consumer's risk \(0.265274\), ASN at AQL \(10.955041\),
and ASN at RQL \(16.233221\). It reduces average inspection at the cost
of accepting more bad lots in this example; it is not a universal improvement.
ATI at AQL is \(17.313145\) and AOQ \(0.009654\), under the solver's
rectifying convention. For double plans, QSB replaces \(n\) with ASN in
the reported ATI and AOQ formulas. This is an aggregate convention: it does
not separately calculate the inspected size conditional on each acceptance
or rejection route.

The CLI can design the smallest \emph{single} plan satisfying binomial bounds:
\begin{lstlisting}
.build/debug/qsb design-acceptance 500 0.01 0.15 0.05 0.10
\end{lstlisting}
The result is \(n=25,c=1\), producer's risk \(0.025759\), and consumer's
risk \(0.093071\). Argument order is lot size, AQL, RQL, producer's risk,
and consumer's risk. The command searches for a plan, rather than evaluating
one already fixed. It does not design double plans, optimize expected cost,
implement nonzero classification errors, or claim conformity with a specific
sampling standard.

In macOS, the form distinguishes single/double plans and their decision
thresholds; it preserves shared quality levels, risks, and economic data when
the type changes. Validate and solve to see the OC curve, AQL/RQL markers,
actual versus nominal risks, and ASN/ATI/AOQ. Economic fields are preserved
in the model; core output contains quality and inspection measures, rather
than an economic optimum. JSON allows the exact plan used to be audited.

\section{Exercises}
\begin{enumerate}
\item Add a count of twenty to the c example and observe how the center and
limits change when it is included in estimation. Discuss when to establish a
reference phase.
\item Compare p charts with sizes of one hundred and one thousand, keeping
proportions fixed. Explain why the same deviation can produce different signals.
\item Design a single plan using the example's bounds and compare ASN with
the \(20,1\) plan. Check actual risks, rather than only nominal ones.
\end{enumerate}


\appendix
\chapter{Terminology reference}
\label{chap:terminos}
\begin{description}
 \item[Backend] Implementation selected to validate or solve a model family.
 The validation option does not compute a solution.
 \item[Balance] Equality accounting for inflows, outflows, and changes in state,
 such as ending inventory equaling initial inventory plus production minus demand.
 \item[Capacity] Resource limit, with a defined unit and time horizon.
 \item[Constraint] Condition limiting the admissible decisions.
 \item[Decision variable] Controllable quantity chosen by the decision maker in the model.
 \item[Diagnostic] Validation message with severity, code, and, when available,
 the path of the affected field.
 \item[Exactness] Characterization of the procedure and its guarantees for the
 supported model, distinct from the number of decimal places displayed.
 \item[Expected value] Probability-weighted average; it does not guarantee that
 payoff in any particular realization.
 \item[Feasibility] Satisfaction of the model's constraints, bounds, and
 conditions; it does not imply optimality.
 \item[Feasible region] Set of solutions satisfying the constraints.
 \item[Global optimum] Feasible solution with the best objective throughout the
 feasible region; a local optimum is compared only with nearby solutions.
 \item[Heuristic] Procedure seeking a good solution without necessarily
 guaranteeing a global optimum.
 \item[Likelihood] Probability of observing evidence conditional on a state;
 it is not the posterior probability of that state.
 \item[Linear relaxation] Model replacing integrality requirements with
 continuous variables; its solution may be inadmissible for the integer problem.
 \item[Makespan] Completion time of the last job in a schedule.
 \item[Normalized JSON] Structured representation of a model or result that
 allows data to be exchanged, inspected, and reproduced across QSB clients.
 \item[Objective function] Expression quantifying the criterion to maximize or minimize.
 \item[Optimal solution] Feasible solution achieving the best objective of the
 problem considered, within the method's guarantees and tolerances.
 \item[Parameter] Data item treated as known by the model during a run.
 \item[Policy] Rule or set of decisions by stage and state.
 \item[Posterior] Probability of a state after observing evidence, calculated
 from priors and likelihoods.
 \item[Seed] Value initializing the pseudorandom generator; it enables a
 simulation to be repeated with the same model, configuration, and implementation.
 \item[Sensitivity] Analysis of how the solution changes when data or assumptions change.
 \item[Shadow price] Marginal change in the objective when a resource varies
 within the range and under the conditions where the dual interpretation is valid.
 \item[Slack] Remaining capacity or margin of a constraint for a solution.
 \item[State] Summary of sufficient information for future decisions in a
 dynamic model; it may also refer to a random condition of the environment.
 \item[Tolerance] Numerical threshold used to compare quantities or accept
 approximations; it must be interpreted alongside the scale of the problem.
 \item[Trace] Record of stages, states, and decisions helping to examine a
 computed policy; it is not necessarily the algorithm's entire internal table.
 \item[Utilization] Fraction of service capacity occupied on average.
 \item[Value of information] Expected improvement from allowing decisions based
 on evidence; EVSI refers to sample information and EVPI to perfect information.
\end{description}


\backmatter
\bibliographystyle{apalike}
\bibliography{references}

\end{document}
